\chapter{Sequences, Series and Limits}\label{ch:sequences}

\begin{center}\itshape
$1-\frac12+\frac13-\frac14+\cdots=\ln2$. Rearrange the same terms as $1-\frac12-\frac14+\frac13-\frac16-\frac18+\cdots$\\
and the sum becomes $\frac12\ln2$. Infinite sums are not as innocent as finite ones.
\end{center}

Write the first six partial sums of $1-\frac12+\frac13-\frac14+\frac15-\frac16$:
\[
1,\quad \frac12,\quad \frac56,\quad \frac7{12},\quad \frac{47}{60},\quad \frac{37}{60}.
\]
As decimals they are $1$, $0.5$, $0.833$, $0.583$, $0.783$, $0.617$. They hop from one side of $\ln2\approx0.693$ to the other, and the hop gets smaller.

Now list the same fractions with one positive term and then two negative ones:
\[
1-\frac12-\frac14+\frac13-\frac16-\frac18+\frac15-\frac1{10}-\frac1{12}+\cdots.
\]
After each completed block of three terms the sums are
\[
0.25,\quad 0.292,\quad 0.308,\quad 0.317.
\]
They are climbing toward $\frac12\ln2\approx0.347$, not toward $\ln2$. A finite sum is deaf to order. These two infinite sums are not.

The previous chapter stopped on a question of the same kind. The derangement proportions $D_n/n!$ sit still in Euler's table once the pack is long. A table is not a definition of settling. This chapter gives the definition, and it separates the sums that may be rearranged from the sums that may not.

\portraits{oresme}{Nicole Oresme (c.\,1320--1382), miniature from a manuscript of his \emph{Trait\'e de la sph\`ere}, c.\,1400--1420. Public domain, via Wikimedia Commons.}{bolzano}{Bernard Bolzano (1781--1848). Public domain, via Wikimedia Commons.}

Around 450~BC, Zeno of Elea argued that motion is impossible. To cross a room you must first cross half of it, then half of what remains, then half of that, and so on: infinitely many tasks. Aristotle reports the argument in the \emph{Physics}, Book~VI, Chapter~9, and answers that the half-rooms take half-times which add to a finite time. He did not have a definition of the sum $\frac12+\frac14+\frac18+\cdots$.

Nicole Oresme, at the University of Paris around 1350, had two infinite sums in his questions on Euclid's geometry. One of them is finite: $\frac12+\frac24+\frac38+\frac4{16}+\cdots=2$. The other, the harmonic series $1+\frac12+\frac13+\cdots$, grows past every bound. The grouping he used is Exercise~\ref{ex:oresme}, and Figure~\ref{fig:cobweb} draws the blocks. In 1703 Guido Grandi treated $1-1+1-1+\cdots$ as equal to $\frac12$, because $1/(1+1)$ expands to that series. Leibniz, writing to Christian Wolff in 1713, defended the same value: a partial sum is $1$ or $0$ according as the number of terms is odd or even, and the two cases should carry equal weight, so the value to take is their average.

The nineteenth century wrote the definitions those arguments needed. Bernard Bolzano, in 1817, proved the intermediate-value theorem from the least-upper-bound property, by the bisection used below. Cauchy's \emph{Cours d'analyse} of 1821, written for his students at the \'Ecole polytechnique, defined convergence. Weierstrass, in his Berlin lectures, fixed the order of the quantifiers in that definition: for every tolerance there exists a rank. Riemann, in the habilitation of 1854 published in 1867, showed that a series which converges by cancellation can be rearranged to any sum you name, or to no sum at all. The epigraph is the smallest case. Exercise~\ref{ex:riemannrearr} is the general fact.

Order, in this chapter, is a sequence that has settled, or a sum whose value does not depend on the list. Chance, in the loose sense Grandi and Leibniz were reaching for, is an average of sums that refuse to settle. Absolute convergence is what makes order irrelevant; averaging is what rescues a sum that never settles.

\section{What a limit is}

For a real number $x$, set $|x|=\max\{x,-x\}$. Then $-|x|\le x\le|x|$, and $|x|\le r$ means $-r\le x\le r$ whenever $r\ge0$. The next two inequalities are the ones every later estimate quotes.

\begin{proposition}[real triangle inequalities]\label{prop:real-triangle}
For real $a,b$,
\[
|a+b|\le|a|+|b|,\qquad \bigl||a|-|b|\bigr|\le|a-b|.
\]
\end{proposition}
\begin{proof}
Add $-|a|\le a\le|a|$ and $-|b|\le b\le|b|$. You get $-(|a|+|b|)\le a+b\le|a|+|b|$, which is the first inequality.

For the second, write $a=(a-b)+b$. The first inequality gives $|a|\le|a-b|+|b|$, so $|a|-|b|\le|a-b|$. Swap $a$ and $b$: $|b|-|a|\le|a-b|$. A number and its negative are both at most $|a-b|$, so the larger of them, which is $\bigl||a|-|b|\bigr|$, is at most $|a-b|$.
\end{proof}

Check $a=3$ and $b=-5$. Then $|a+b|=2$ and $|a|+|b|=8$. Also $\bigl||3|-|{-5}|\bigr|=2$ and $|3-(-5)|=8$, so $2\le8$. Equality in the second inequality is easy to hit: $|5|-|3|=2=|5-3|$.

\begin{reminder}[what a limit is]
A sequence $(u_n)_{n\ge0}$ of real numbers \emph{converges} to $\ell\in\R$ if
\[
\forall\varepsilon>0,\ \exists N\in\N,\ \forall n\ge N,\quad|u_n-\ell|\le\varepsilon.
\]
We write $u_n\to\ell$, or $\lim u_n=\ell$. Read it as a game. Your opponent names a tolerance $\varepsilon>0$, as small as she likes. You must answer with a rank $N$ beyond which \emph{every} term lies within $\varepsilon$ of $\ell$. You win if you can always answer.

Try $u_n=\frac1n$ for $n\ge1$, and the claim $\ell=0$. If she names $\varepsilon=\frac1{10}$, any $N>10$ works: for $n\ge11$, $0<\frac1n\le\frac1{11}<\frac1{10}$. If she names a general $\varepsilon>0$, any natural number $N>\frac1\varepsilon$ works, once we know such an $N$ exists. That existence is the Archimedean property, proved in the next section from completeness. The game itself is only the inequality.

The sequence $(-1)^n$ converges to nothing. For $\varepsilon=\frac12$, no rank works: consecutive terms are $2$ apart, so they cannot both lie in an interval of length $1$.

A sequence \emph{tends to $+\infty$} if for every real $A$ there is an $N$ with $u_n\ge A$ for all $n\ge N$. \emph{Diverges} means that the sequence does not converge to a real number. Tending to $\pm\infty$, and oscillating, are both ways to diverge.
\end{reminder}

Negating the definition is a check that the quantifiers are in the right order. The sentence $u_n\not\to\ell$ means that some $\varepsilon>0$ defeats every proposed $N$: for every $N$ there is an $n\ge N$ with $|u_n-\ell|>\varepsilon$. One early bad term does not spoil a limit. The bad index is allowed to run off to infinity as $N$ grows.

\portraits{cauchy}{Augustin-Louis Cauchy (1789--1857). Illustration from \emph{Das neunzehnte Jahrhundert in Bildnissen}, Berlin, 1901. Public domain, via Wikimedia Commons.}{weierstrass}{Karl Weierstrass (1815--1897). Portrait by Conrad Fehr, 1895. Public domain, via Wikimedia Commons.}

\begin{figure}[ht]
\centering
\fig{ch05-epsilon}
\caption{Left: the game of $\varepsilon$ and $N$ for $u_n=1+\frac{(-1)^n}{n}$, which converges to $1$. For $\varepsilon=0.15$ the band around $1$ contains every term from $N=7$ on. The first six terms, drawn in rust, are allowed to misbehave. A smaller $\varepsilon$ asks for a larger $N$, and that is the whole definition. Right: peaks of a bounded sequence, drawn in rust. A peak is at least as large as every later term. Infinitely many peaks give a non-increasing subsequence. Finitely many of them leave you room to climb.}
\label{fig:epsilon}
\end{figure}

On the left of Figure~\ref{fig:epsilon}, take $\varepsilon=0.15$ and $\ell=1$. The term with $n=6$ is $1+\frac16\approx1.167$, which is still outside the band. The term with $n=7$ is $1-\frac17\approx0.857$, and $\bigl|0.857-1\bigr|\approx0.143\le0.15$. For every later $n$, $\bigl|\frac{(-1)^n}n\bigr|\le\frac17\le0.15$. So $N=7$ answers this particular $\varepsilon$. It does not answer $\varepsilon=0.01$. That is not a defect. The definition asks for an $N$ after $\varepsilon$ has been named, not for one $N$ that works for every $\varepsilon$ at once.

\section{Completeness}

Every argument below that produces a limit, rather than computing one already in hand, uses one axiom.

\begin{reminder}[completeness of $\R$]
Every non-empty subset of $\R$ that is bounded above has a least upper bound, called its \emph{supremum}.
\end{reminder}

The natural numbers are not bounded above in $\R$. If they were, $M=\sup\N$ would exist. Some $n\in\N$ would then satisfy $n>M-1$, because $M-1$ is not an upper bound. But then $n+1>M$, and $n+1$ is still a natural number, which contradicts the choice of $M$. So no such $M$ exists: no real number is an upper bound for $\N$. For every real $x$ there is a natural number $n>x$. This is the Archimedean property, and it is why the rank $N>\frac1\varepsilon$ in the game for $\frac1n$ can be chosen.

The same fact cuts the line into unit steps, which is the floor the counting chapter was waiting for.

\begin{proposition}[the floor]\label{prop:floor}
Every real number $x$ has a unique integer floor $\lfloor x\rfloor$, characterised by
\[
\lfloor x\rfloor\le x<\lfloor x\rfloor+1.
\]
\end{proposition}
\begin{proof}
\emph{Existence.} By the Archimedean property there is a natural number $n>-x$, so the integer $-n$ satisfies $-n<x$, and the set
\[
S=\{m\in\Z:m\le x\}
\]
is non-empty. It is bounded above by $x$, so $s=\sup S$ exists in $\R$. Some integer $m\in S$ satisfies $m>s-1$, because $s-1$ is not an upper bound of $S$. Then $m\le x$ and $m+1>s$. The integer $m+1$ cannot lie in $S$, so $m+1>x$. Thus $m\le x<m+1$.

\emph{Uniqueness.} Suppose $m$ and $m'$ are integers with $m\le x<m+1$ and $m'\le x<m'+1$. If $m<m'$, then $m+1\le m'$, so $x<m+1\le m'\le x$, which is absurd. The symmetric argument rules out $m'>m$. So $m=m'$.
\end{proof}

Check the endpoint and the negative side, which is the one that surprises people. $\lfloor 2\rfloor=2$, because $2\le2<3$. $\lfloor2.3\rfloor=2$, because $2\le2.3<3$. $\lfloor-1.2\rfloor=-2$, because $-2\le-1.2<-1$. The integer $-1$ is too large: $-1\not\le-1.2$. Figure~\ref{fig:floor} is those two points. The filled dot is included in the step; the open dot is the next integer, which is not.

\begin{figure}[ht]
\centering
\fig{ch05-floor}
\caption{The floor is the left end of the unit step that contains $x$. The step for $2.3$ is $[2,3)$, so $\lfloor2.3\rfloor=2$. The step for $-1.2$ is $[-2,-1)$, so $\lfloor-1.2\rfloor=-2$. The negative case steps down, not toward zero.}
\label{fig:floor}
\end{figure}

Completeness is also what $\Q$ does not have. The next proposition is the usual witness. It uses Euclid's lemma from Chapter~\ref{ch:sums}: a prime that divides a product divides one of the factors.

\begin{proposition}[$\sqrt2$ is not rational]\label{prop:sqrt2}
There is no rational number $q$ with $q^2=2$.
\end{proposition}
\begin{proof}
Suppose $q=a/b$ with $a,b$ positive integers and $\gcd(a,b)=1$. Then $a^2=2b^2$, so $2$ divides $a^2=a\cdot a$. Since $2$ is prime, $2$ divides $a$. Write $a=2c$. Then $4c^2=2b^2$, so $b^2=2c^2$, and the same lemma says that $2$ divides $b$. Thus $2$ divides both $a$ and $b$, which contradicts $\gcd(a,b)=1$.
\end{proof}

Squaring is strictly increasing on the positive reals: if $0\le a<b$, then $b^2-a^2=(b-a)(a+b)>0$.

Let $E=\{q\in\Q:q>0,\ q^2<2\}$. The set is non-empty, since $1\in E$, and it is bounded above in $\Q$, since every element is $<2$. It has no least upper bound in $\Q$. To see this, let $r$ be a positive rational upper bound. Then $r^2\ne2$ by Proposition~\ref{prop:sqrt2}. The case $r^2<2$ is also impossible. Choose a natural number $n>(2r+1)/(2-r^2)$, using the Archimedean property. For every integer $n\ge1$ one has $\frac1{n^2}\le\frac1n$, so
\[
\Bigl(r+\frac1n\Bigr)^2=r^2+\frac{2r}n+\frac1{n^2}\le r^2+\frac{2r+1}n<2,
\]
so $r+\frac1n$ lies in $E$ and is larger than $r$. Thus any positive rational upper bound $r$ satisfies $r^2>2$. For such an $r$, the rational
\[
r'=\frac12\Bigl(r+\frac2r\Bigr)
\]
is positive and strictly smaller than $r$, because $r-r'=(r^2-2)/(2r)>0$. Also $(r')^2-2=(r^2-2)^2/(4r^2)>0$, so $(r')^2>2$. This $r'$ is still an upper bound of $E$. If some $q\in E$ satisfied $q>r'$, then $q>0$ and squaring would give $q^2>(r')^2>2$, which contradicts $q\in E$. So $r'$ is a strictly smaller positive rational upper bound. No upper bound in $\Q$ is the least one. This is the sense in which $\Q$ is not complete.

The same constructions inside $\R$ produce the root. Let $F=\{x\in\R:x>0,\ x^2<2\}$. The set is non-empty and bounded above, so $\sigma=\sup F$ exists. If $\sigma^2<2$, choose an integer $n>(2\sigma+1)/(2-\sigma^2)$. The estimate written above, with $\sigma$ in place of $r$, gives $\bigl(\sigma+\frac1n\bigr)^2<2$, so $\sigma+\frac1n$ lies in $F$ and exceeds $\sigma$. If $\sigma^2>2$, then $\sigma'=\frac12\bigl(\sigma+\frac2\sigma\bigr)$ satisfies $0<\sigma'<\sigma$ and $(\sigma')^2-2=(\sigma^2-2)^2/(4\sigma^2)>0$. Any $x>\sigma'$ has $x^2>(\sigma')^2>2$, so $\sigma'$ is a smaller upper bound of $F$. Both alternatives contradict the choice of $\sigma$. Therefore $\sigma^2=2$. This positive real is written $\sqrt2$. It is the only one: squaring is increasing on the positive reals. Newton's sequence below converges to it.

\begin{propositionx}[first properties of limits]
Let $u_n\to\ell$ and $v_n\to m$, with $\ell,m\in\R$.
\begin{enumerate}[label=(\alph*)]
\item The limit is unique, and a convergent sequence is bounded.
\item If $u_n\le v_n$ for all $n$ large, then $\ell\le m$. A strict inequality may die in the limit: $\frac1n>0$, but $\lim\frac1n=0$.
\item \textbf{Squeeze.} If $u_n\le w_n\le v_n$ for all $n$ large and $\ell=m$, then $w_n\to\ell$.
\item $u_n+v_n\to\ell+m$ and $u_nv_n\to\ell m$. If $m\ne0$, then $v_n\ne0$ for $n$ large and $\frac{u_n}{v_n}\to\frac\ell m$.
\item \textbf{Subsequences.} If $n_0<n_1<n_2<\cdots$, then $u_{n_k}\to\ell$ as $k\to\infty$.
\end{enumerate}
\end{propositionx}
\begin{proof}
\emph{Uniqueness and boundedness.} Suppose also $u_n\to\ell'$ with $\ell\ne\ell'$. Set $\varepsilon=\frac{|\ell-\ell'|}3$. For $n$ large, $|u_n-\ell|\le\varepsilon$ and $|u_n-\ell'|\le\varepsilon$, so
\[
|\ell-\ell'|\le|\ell-u_n|+|u_n-\ell'|\le2\varepsilon=\frac{2}{3}|\ell-\ell'|,
\]
which is impossible. For boundedness, take $\varepsilon=1$. Past some $N$, every term lies in $[\ell-1,\ell+1]$. The finitely many terms before $N$ have some maximum absolute value $M$. Every term then satisfies $|u_n|\le\max\bigl(M,|\ell|+1\bigr)$.

\emph{Inequalities.} Suppose $\ell>m$, and set $\varepsilon=\frac{\ell-m}3$. For $n$ large, $u_n\ge\ell-\varepsilon>m+\varepsilon\ge v_n$, which contradicts $u_n\le v_n$. The strict inequality $\frac1n>0$ has limit $0$, so a strict inequality need not stay strict.

\emph{Squeeze.} Given $\varepsilon>0$, both $u_n$ and $v_n$ lie in $[\ell-\varepsilon,\ell+\varepsilon]$ for $n$ large. Any $w_n$ between them lies in the same interval.

\emph{Sum.} $|u_n+v_n-\ell-m|\le|u_n-\ell|+|v_n-m|$. Given $\varepsilon>0$, choose $N$ so that for $n\ge N$ each of those two terms is at most $\frac\varepsilon2$. The sum is then at most $\varepsilon$.

\emph{Product.} Write $u_nv_n-\ell m=(u_n-\ell)v_n+\ell(v_n-m)$. By the boundedness already proved, $|v_n|\le B$ for some $B$ and every $n$. Then
\[
|u_nv_n-\ell m|\le B|u_n-\ell|+|\ell|\,|v_n-m|.
\]
Given $\varepsilon>0$, make $|u_n-\ell|\le\frac{\varepsilon}{2(B+1)}$ and $|v_n-m|\le\frac{\varepsilon}{2(|\ell|+1)}$. Each piece is then at most $\frac\varepsilon2$. The extra $1$ in the denominators is only there so that we never divide by zero.

\emph{Quotient.} By the product rule it is enough to show $\frac1{v_n}\to\frac1m$. Take $\varepsilon=\frac{|m|}2$ in the definition of $v_n\to m$. For $n$ large, $|v_n-m|\le\frac{|m|}2$, so $|v_n|\ge\frac{|m|}2>0$. Then
\[
\Bigl|\frac1{v_n}-\frac1m\Bigr|=\frac{|m-v_n|}{|v_n|\,|m|}\le\frac{2|v_n-m|}{m^2}.
\]
Given a tolerance $\eta>0$, make $|v_n-m|\le\eta\, m^2/2$. The displayed quantity is then at most $\eta$.

\emph{Subsequences.} The integers satisfy $n_k\ge k$. This holds for $k=0$ because $n_0\ge0$, and if it holds for $k$ then $n_{k+1}\ge n_k+1\ge k+1$. Once $k\ge N$, every $n_k$ is at least $N$, so $|u_{n_k}-\ell|\le\varepsilon$.
\end{proof}

\begin{reminder}[the notations $o$, $O$ and $\sim$]
Let $v_n\ne0$ for all $n$ large.
\begin{itemize}
\item $u_n=o(v_n)$ means $\frac{u_n}{v_n}\to0$. Thus $o(1)$ is any sequence tending to $0$, and $n=o(n^2)$.
\item $u_n=O(v_n)$ means $\bigl|\frac{u_n}{v_n}\bigr|\le C$ for some constant $C$ and all $n$ large. Thus $O(1)$ is any bounded sequence, and $\sin n=O(1)$.
\item $u_n\sim v_n$ means $\frac{u_n}{v_n}\to1$, that is $u_n=v_n+o(v_n)$. Thus $n^2+n\sim n^2$.
\end{itemize}
Equivalents may be multiplied and divided. They are not a licence to add or to subtract. At $n=100$,
\[
\frac{n+1}n=1.01,\qquad\frac{-n+\sqrt n}{-n}=0.9,
\]
so both equivalences $n+1\sim n$ and $-n+\sqrt n\sim-n$ are already visible. The left sides add to $1+\sqrt n=11$. The right sides add to $0$. Eleven is not close to zero. Subtraction fails in the same way. At $n=100$, both $n+1\sim n$ and $n\sim n$ hold, and the left sides differ by $1$ while the right sides differ by $0$. Keep the first term that does not cancel, or keep a remainder you can bound.

The same symbols are used for functions, with $x\to a$ in place of $n\to\infty$. For instance $\sin x\sim x$ as $x\to0$: the admitted derivative $\sin'(0)=\cos0=1$ says $\frac{\sin x}x\to1$. Read $o$ and $O$ from left to right only. The sentence $u_n=O(v_n)$ talks about $u_n$. It is not an equation you may rearrange.
\end{reminder}

\begin{examplex}[the geometric series]
For $|q|<1$,
\[
\sum_{n=0}^\infty q^n=\lim_{N\to\infty}\frac{1-q^{N+1}}{1-q}=\frac1{1-q}.
\]
\end{examplex}

The finite sum is the usual one: $(1-q)\sum_{n=0}^N q^n=1-q^{N+1}$, and $q\ne1$ because $|q|<1$. The limit is $\frac1{1-q}$ once $q^{N+1}\to0$. If $q=0$ that is immediate. If $0<|q|<1$, set $\frac1{|q|}=1+h$ with $h>0$. Bernoulli's inequality, Proposition~\ref{prop:bernoulli}, gives $(1+h)^N\ge1+Nh$, which tends to infinity. Therefore $|q|^N=1/(1+h)^N\to0$.

Check $q=\frac12$ and $N=3$. The partial sum is $1+\frac12+\frac14+\frac18=\frac{15}8$. The closed form gives $\frac{1-(1/2)^4}{1/2}=2\bigl(1-\frac1{16}\bigr)=\frac{15}8$. The infinite sum is $2$.

Zeno's half-rooms are the tail of this series. Figure~\ref{fig:zeno} fills a segment of length $1$ by pieces of length $\frac12$, $\frac14$, $\frac18$, $\frac1{16}$, and a remainder. The partial sum after $N$ pieces is $1-2^{-N}$, which tends to $1$. So $\sum_{n\ge1}2^{-n}=1$. There are infinitely many pieces and a finite length. The sum is a limit of finite sums, not a task that has to be finished by hand.

\begin{figure}[ht]
\centering
\fig{ch05-zeno}
\caption{A segment of length $1$, cut into $\frac12+\frac14+\frac18+\frac1{16}$ and a grey tail of length $\frac1{16}$. The labels above the cuts are the partial sums $\frac12$, $\frac34$, $\frac78$ and $\frac{15}{16}$. The tail is exactly what is still missing from $1$.}
\label{fig:zeno}
\end{figure}

Oresme's other finite sum is the same geometric series with weights. Let $S_N=\sum_{k=1}^N k/2^k$. Then
\[
\frac12 S_N=\sum_{k=1}^N\frac{k}{2^{k+1}}=\sum_{k=2}^{N+1}\frac{k-1}{2^k},
\]
and subtracting gives
\[
S_N-\frac12 S_N=\frac12+\sum_{k=2}^N\frac1{2^k}-\frac{N}{2^{N+1}}=\sum_{k=1}^N\frac1{2^k}-\frac{N}{2^{N+1}}=1-2^{-N}-\frac{N}{2^{N+1}}.
\]
Hence $S_N=2-(N+2)2^{-N}$. For $N\ge2$, the binomial theorem supplies $2^N=(1+1)^N\ge\binom N2=N(N-1)/2$, so
\[
0<\frac{N}{2^N}\le\frac{2}{N-1}\to0.
\]
The remaining piece $2^{1-N}$ tends to $0$ as well, so $(N+2)2^{-N}\to0$. Therefore $S_N\to2$. The first values are $\frac12$, $1$, $\frac{11}8$, $\frac{13}8$, already $1.625$ at $N=4$, and $1.988$ at $N=10$. Infinitely many positive terms, and the sum is $2$.

Grandi's series is the series that does not get this courtesy. Its terms are $1,-1,1,-1,\ldots$, which do not tend to $0$. If the partial sums $S_n$ tended to a finite $S$, the terms $u_n=S_n-S_{n-1}$ would tend to $S-S=0$. So $1-1+1-1+\cdots$ has no sum in the sense of this chapter. The partial sums are $1,0,1,0,\ldots$. Leibniz's equal weight is an average of those partial sums, drawn in Figure~\ref{fig:grandi}. The averages really do tend to $\frac12$. That fact is Ces\`aro's theorem, applied to a sequence which itself refuses to converge. It is not a value of the series.

\begin{figure}[ht]
\centering
\fig{ch05-grandi}
\caption{Partial sums of $1-1+1-1+\cdots$, in charcoal, alternate between $1$ and $0$. The running averages, in rust, are $1,\frac12,\frac23,\frac12,\frac35,\frac12,\ldots$ and they settle at $\frac12$. The average has a limit. The series does not.}
\label{fig:grandi}
\end{figure}

\begin{think}
Must a sequence with $u_{n+1}-u_n\to0$ converge? Try $u_n=\sqrt n$, then $u_n=\ln n$. What would you have to ask of the gaps $u_p-u_q$, for all large $p$ and $q$, before convergence was forced?
\end{think}

\section{Sequences that have to converge}

The four tests in the next proposition are ways to know that a limit exists without already knowing its name. The geometric series did not need them: the partial sums were computed. Most of the limits in the rest of this volume are not like that.

\begin{proposition}[convergence criteria]\label{prop:seqcrit}
\begin{enumerate}[label=(\alph*)]
\item A non-decreasing sequence converges if and only if it is bounded above, and then it converges to its supremum. A non-increasing sequence converges if and only if it is bounded below, and then it converges to its infimum.
\item \textbf{Adjacent sequences.} If $(u_n)$ is non-decreasing, $(v_n)$ is non-increasing, and $v_n-u_n\to0$, then both sequences converge to the same limit $\ell$, and $u_n\le\ell\le v_n$.
\item \textbf{Bolzano--Weierstrass.} Every bounded sequence has a convergent subsequence.
\item \textbf{Cauchy.} The sequence $(u_n)$ converges if and only if for every $\varepsilon>0$ there is an $N$ such that $|u_p-u_q|\le\varepsilon$ whenever $p,q\ge N$.
\end{enumerate}
\end{proposition}
\begin{proof}
\emph{Monotone sequences.} A convergent sequence is bounded, by the previous proposition. Conversely, let $(u_n)$ be non-decreasing and bounded above, and set $\ell=\sup_n u_n$, which exists by completeness. Given $\varepsilon>0$, the number $\ell-\varepsilon$ is not an upper bound, so $u_N>\ell-\varepsilon$ for some $N$. For every $n\ge N$, monotonicity gives $\ell-\varepsilon<u_N\le u_n\le\ell$. That is $u_n\to\ell$. If instead $(u_n)$ is non-increasing and bounded below, apply this case to $(-u_n)$. The supremum of $(-u_n)$ is $-\inf_n u_n$.

\emph{Adjacent sequences.} The difference $v_n-u_n$ is non-increasing, because $v$ falls and $u$ rises. A non-increasing sequence that tends to $0$ cannot have a negative term: if $v_k-u_k<0$, then every later difference is $\le v_k-u_k<0$, so the limit would be $\le v_k-u_k$ and could not be $0$. Thus $v_n\ge u_n$ for every $n$, and $u_0\le u_n\le v_n\le v_0$. Part (a) gives $u_n\to\ell$ and $v_n\to\ell'$. Then $\ell'-\ell=\lim(v_n-u_n)=0$. Monotonicity keeps $u_n\le\ell$ and $v_n\ge\ell'$.

\emph{Bolzano--Weierstrass.} Call $n$ a \emph{peak} if $u_n\ge u_m$ for every $m>n$. The right half of Figure~\ref{fig:epsilon} marks six peaks. If there are infinitely many peaks $n_1<n_2<\cdots$, the subsequence $(u_{n_k})$ is non-increasing and bounded, so it converges by (a). If there are only finitely many peaks, start after the last one. The index you are on is not a peak, so some later term is strictly larger. Repeating this builds a strictly increasing subsequence. It is bounded, so it converges by (a).

\emph{Cauchy's criterion.} If $u_n\to\ell$, then $|u_p-u_q|\le|u_p-\ell|+|\ell-u_q|$, and both pieces are small once $p$ and $q$ are large. Conversely, a Cauchy sequence is bounded: take $\varepsilon=1$, and past the corresponding $N$ every term is within $1$ of $u_N$. Part (c) gives a subsequence $u_{n_k}\to\ell$. Now fix $\varepsilon>0$ and choose $N$ so that $|u_p-u_q|\le\frac\varepsilon2$ for all $p,q\ge N$. Choose $k$ with $n_k\ge N$ and $|u_{n_k}-\ell|\le\frac\varepsilon2$. For every $p\ge N$,
\[
|u_p-\ell|\le|u_p-u_{n_k}|+|u_{n_k}-\ell|\le\varepsilon.
\]
So $u_p\to\ell$.
\end{proof}

Figure~\ref{fig:squeeze} is the adjacent case: once the gap tends to $0$, the two graphs have nowhere to go except a common height. Cauchy's criterion is the form you use when you cannot see the limit. The sequences $\sqrt n$ and $\ln n$ from the box satisfy $u_{n+1}-u_n\to0$, but $u_{2n}-u_n\not\to0$. For the square root, $\sqrt{2n}-\sqrt n=\sqrt n(\sqrt2-1)$, which tends to infinity. Consecutive gaps getting small is not the Cauchy condition. The Cauchy condition asks that \emph{every} pair of late terms be close, including terms that are very far apart in the index.

\begin{figure}[ht]
\centering
\fig{ch05-squeeze}
\caption{Adjacent sequences. The lower sequence rises, the upper one falls, and the gap tends to $0$. Completeness supplies one height $\ell$ between them. Neither graph can cross $\ell$.}
\label{fig:squeeze}
\end{figure}

\begin{examplex}[Newton's sequence for $\sqrt2$]
Let $u_0=1$ and $u_{n+1}=\frac12\bigl(u_n+\frac2{u_n}\bigr)$. Then $u_n\to\sqrt2$.
\end{examplex}

The first terms are $u_1=\frac32$ and $u_2=\frac{17}{12}=1.4166\ldots$. The next term already lies within $3\times10^{-6}$ of $\sqrt2\approx1.41421356$. Exercise~\ref{ex:newton} explains the speed. Here we only want the limit.

All terms are positive. For $n\ge0$,
\[
u_{n+1}-\sqrt2=\frac{u_n^2+2-2u_n\sqrt2}{2u_n}=\frac{(u_n-\sqrt2)^2}{2u_n}\ge0,
\]
so $u_n\ge\sqrt2$ once $n\ge1$. Then
\[
u_{n+1}-u_n=\frac{2-u_n^2}{2u_n}\le0
\]
for $n\ge1$. From index $1$ on, the sequence is non-increasing and bounded below by $\sqrt2$, so it converges to some $\ell\ge\sqrt2$. Pass to the limit in the recurrence, using the sum, product and quotient rules:
\[
\ell=\frac\ell2+\frac1\ell.
\]
Multiply by $2\ell$, using $\ell>0$: $2\ell^2=\ell^2+2$, so $\ell^2=2$, and $\ell=\sqrt2$.

The pattern is worth keeping. If $u_{n+1}=f(u_n)$ and $f$ is continuous, a limit of the sequence has to be a fixed point of $f$. Monotonicity, read off the sign of $f(x)-x$, decides whether the sequence converges. Figure~\ref{fig:cobweb} draws the same idea for $u_{n+1}=\sin u_n$, which is Exercise~\ref{ex:sin}.

\begin{figure}[ht]
\centering
\fig{ch05-cobweb}
\caption{Left: the cobweb of $u_{n+1}=\sin u_n$ from $u_0=1.5$. A vertical segment runs from the diagonal up to $y=\sin x$; a horizontal segment carries the new value back to the diagonal. The steps crawl toward the fixed point $0$. Right: Oresme's blocks for the harmonic series. After the first bar, group $1$, then $2$, then $4$, then $8$ bars. Each of those blocks has area at least $\frac12$, so the partial sums are unbounded.}
\label{fig:cobweb}
\end{figure}

One block is enough to see the picture. The four terms $\frac15+\frac16+\frac17+\frac18$ are the bars from just after $4$ up to $8$. Each is at least $\frac18$, and there are four of them, so the block is at least $\frac12$. Exercise~\ref{ex:oresme} does this for every block of length $2^{j-1}$.

\begin{think}
If $u_n\to\ell$, the averages $\frac{u_1+\cdots+u_n}n$ ought to tend to $\ell$ as well: a fixed number of early terms is drowned by the denominator. Is the converse true? The partial sums in Figure~\ref{fig:grandi} are a test case.
\end{think}

\begin{theorem}[Ces\`aro]\label{thm:cesaro}
If $u_n\to\ell\in\R$, then $\dfrac{u_1+\cdots+u_n}n\to\ell$.
\end{theorem}
\begin{proof}
Replacing $u_n$ by $u_n-\ell$, we may assume $\ell=0$. The averages shift by the same constant, so this changes nothing about the claim. Let $\varepsilon>0$, and choose $N$ with $|u_n|\le\varepsilon$ for all $n>N$. For $n>N$,
\[
\Bigl|\frac{u_1+\cdots+u_n}n\Bigr|\le\frac{|u_1|+\cdots+|u_N|}n+\frac{n-N}n\,\varepsilon.
\]
The numerator of the first fraction is fixed, so that fraction tends to $0$. For $n$ large it is at most $\varepsilon$, and the second fraction is at most $\varepsilon$. The average is then at most $2\varepsilon$ in absolute value.
\end{proof}

A corollary drops out at once. If $u_{n+1}-u_n\to\ell$, then $\frac{u_n}n\to\ell$. The differences $d_k=u_{k+1}-u_k$ for $k=0,\ldots,n-1$ have average $(u_n-u_0)/n$. Ces\`aro says this average tends to $\ell$, and $u_0/n\to0$, so $u_n/n\to\ell$.

Figure~\ref{fig:grandi} is the converse failing in public. Let $s_n$ be the partial sums of $1-1+1-1+\cdots$, so $s_n=1$ for odd $n$ and $s_n=0$ for even $n$. This sequence does not converge. Its averages do. If $n=2m$, the sum of the first $n$ partial sums is $m\cdot1+m\cdot0=m$, and the average is $\frac12$. If $n=2m+1$, the sum is $m+1$, and the average is $\frac{m+1}{2m+1}$, which tends to $\frac12$. Leibniz's $\frac12$ is this limit of averages. Ces\`aro's theorem says that a genuine limit would have forced the averages to agree with it. The converse is false, and Grandi's series is the counterexample.

Exercise~\ref{ex:stolz} is the Stolz--Ces\`aro theorem, the same idea with a more general denominator. It is the discrete cousin of l'H\^opital's rule, which waits for Chapter~\ref{ch:calculus}.

\section{Continuous functions on a segment}

The preface left five facts about functions on a segment to be proved, not assumed: intermediate values, extreme values, uniform continuity, Rolle's theorem, and the mean value theorem. All five come out of Proposition~\ref{prop:seqcrit}.

\begin{reminder}[continuity and derivatives]
Let $f$ be defined on an interval $I$, and let $a$ be an accumulation point of $I$: every neighbourhood of $a$ contains a point of $I$ other than $a$. We say $f(x)\to L$ as $x\to a$ if for every $\varepsilon>0$ there is a $\delta>0$ such that $|f(x)-L|\le\varepsilon$ whenever $x\in I$ and $0<|x-a|\le\delta$. When $a\in I$, the function $f$ is \emph{continuous at $a$} if $f(x)\to f(a)$ as $x\to a$, and \emph{continuous on $I$} if it is continuous at every point of $I$.

At an interior point $a$, $f$ is \emph{differentiable at $a$} if $\frac{f(x)-f(a)}{x-a}$ has a limit as $x\to a$. That limit is $f'(a)$. At an endpoint we use a one-sided derivative only when we say so. On a singleton, every function is continuous by convention, and no derivative is defined. A differentiable function is continuous:
\[
f(x)-f(a)=\frac{f(x)-f(a)}{x-a}\,(x-a)\to f'(a)\cdot0=0.
\]

Sequences see a limit. The function $f$ is continuous at $a\in I$ if and only if $f(x_n)\to f(a)$ for every sequence $(x_n)$ in $I$ with $x_n\to a$. If $f$ is continuous at $a$ and $x_n\to a$, take the $\delta$ of continuity and then an $N$ with $|x_n-a|\le\delta$ for $n\ge N$. For those $n$, $|f(x_n)-f(a)|\le\varepsilon$. Conversely, if $f$ is not continuous at $a$, some $\varepsilon>0$ defeats every $\delta$. Taking $\delta=\frac1n$, we find $x_n\in I$ with $|x_n-a|\le\frac1n$ and $|f(x_n)-f(a)|>\varepsilon$. Then $x_n\to a$, but $f(x_n)\not\to f(a)$.

Sums, products and quotients of continuous functions are continuous where the denominator does not vanish, and so are composites. This is the proposition on operations, read along sequences. Polynomials are continuous. So are $\exp$, $\ln$, $\cos$ and $\sin$, whose derivatives the preface admits: $\exp'=\exp$, $\ln'(x)=\frac1x$ for $x>0$, $\sin'=\cos$ and $\cos'=-\sin$.
\end{reminder}

\begin{proposition}[elementary differentiation rules]\label{prop:early-differentiation}
Let $f$ and $g$ be differentiable at an interior point $a$, and let $\lambda,\mu\in\R$.
\begin{enumerate}[label=(\alph*)]
\item $(\lambda f+\mu g)'(a)=\lambda f'(a)+\mu g'(a)$ and $(fg)'(a)=f'(a)g(a)+f(a)g'(a)$. If $g(a)\ne0$, then
\[
\Bigl(\frac fg\Bigr)'(a)=\frac{f'(a)g(a)-f(a)g'(a)}{g(a)^2}.
\]
\item If $h$ is differentiable at $f(a)$ and $h\circ f$ is defined near $a$, then $(h\circ f)'(a)=h'(f(a))\,f'(a)$.
\end{enumerate}
\end{proposition}
\begin{proof}
Linearity is the limit of the difference quotients, using the sum rule for sequences.

For the product, add and subtract $f(a)g(x)$:
\[
\frac{f(x)g(x)-f(a)g(a)}{x-a}=\frac{f(x)-f(a)}{x-a}\,g(x)+f(a)\,\frac{g(x)-g(a)}{x-a}.
\]
Differentiability of $g$ at $a$ implies continuity of $g$ at $a$, so $g(x)\to g(a)$. Each piece has a limit, and the sum of the limits is $f'(a)g(a)+f(a)g'(a)$.

For the quotient, differentiate $\frac1g$ first:
\[
\frac1{x-a}\Bigl(\frac1{g(x)}-\frac1{g(a)}\Bigr)=-\frac{g(x)-g(a)}{x-a}\cdot\frac1{g(x)g(a)}.
\]
The right side tends to $-g'(a)/g(a)^2$, because $g(x)\to g(a)\ne0$. The product rule applied to $f\cdot\frac1g$ gives the quotient formula.

For the chain rule, set $b=f(a)$ and rewrite the definition of $h'(b)$ as
\[
h(y)-h(b)=(y-b)\bigl(h'(b)+\eta(y)\bigr),
\]
where $\eta(y)\to0$ as $y\to b$, and $\eta(b)=0$. This formula also makes sense at $y=b$, which matters because $f(x)$ may equal $f(a)$ for $x$ arbitrarily close to $a$. You cannot divide by $f(x)-f(a)$. Put $y=f(x)$:
\[
\frac{h(f(x))-h(f(a))}{x-a}=\frac{f(x)-f(a)}{x-a}\bigl(h'(b)+\eta(f(x))\bigr).
\]
As $x\to a$, continuity of $f$ gives $f(x)\to b$, so $\eta(f(x))\to0$. The right side tends to $f'(a)h'(b)$.
\end{proof}

Define $x^\alpha=\exp(\alpha\ln x)$ for $x>0$ and real $\alpha$. The chain rule and the admitted derivatives of $\exp$ and $\ln$ give
\[
\frac{\mathrm d}{\mathrm dx}x^\alpha=\exp(\alpha\ln x)\cdot\frac\alpha x=\alpha x^{\alpha-1}.
\]
Check $\alpha=2$: $\exp(2\ln x)\cdot\frac2x=x^2\cdot\frac2x=2x$, which is the derivative you already trust. The comparisons of series later in the chapter use this rule and nothing more elaborate.

\begin{think}
A continuous function on $[0,1]$ with $f(0)=-1$ and $f(1)=1$ has to cross zero, if you draw it. Now restrict $f(x)=x^2-2$ to the rationals in $[0,2]$. It is negative at $0$, positive at $2$, and never zero, by Proposition~\ref{prop:sqrt2}. Which step of the proof below is the one that fails in $\Q$?
\end{think}

\begin{theorem}[intermediate values]\label{thm:ivt}
Let $f$ be continuous on $[a,b]$ with $f(a)<0<f(b)$. Then $f(c)=0$ for some $c\in\,]a,b[$. Consequently, the image of an interval under a continuous function is an interval.
\end{theorem}
\begin{proof}
The argument is bisection, drawn in Figure~\ref{fig:bisection}. Set $a_0=a$ and $b_0=b$. Given $[a_n,b_n]$ with $f(a_n)\le0\le f(b_n)$, let $m$ be the midpoint. If $f(m)\le0$, set $a_{n+1}=m$ and $b_{n+1}=b_n$. Otherwise set $a_{n+1}=a_n$ and $b_{n+1}=m$. Then $f(a_n)\le0\le f(b_n)$ for every $n$. The sequence $(a_n)$ is non-decreasing, $(b_n)$ is non-increasing, and
\[
b_n-a_n=\frac{b-a}{2^n}\to0.
\]
Proposition~\ref{prop:seqcrit}(b) says that both sequences converge to a common limit $c\in[a,b]$. Continuity gives $f(a_n)\to f(c)$ and $f(b_n)\to f(c)$. Weak inequalities survive the limit, so $f(c)\le0\le f(c)$, hence $f(c)=0$. Since $f(a)$ and $f(b)$ are nonzero, $c$ is not an endpoint.

For the image of an interval, suppose $f(u)<v<f(w)$. On the segment joining $u$ and $w$, one endpoint value lies below $v$ and the other lies above it. Apply the theorem to $f-v$ if the left endpoint is the smaller value, and to $v-f$ if the left endpoint is the larger one.

The step that fails over $\Q$ is the appeal to adjacent sequences. Nothing else in the construction used a real number that a rational bisection would not have produced. The midpoints of a rational interval stay rational. Their common limit need not.
\end{proof}

Try the first cut on the function in the figure. Let $f(x)=x^3-2x-1$ on $[1,2.2]$. Then $f(1)=-2<0$ and $f(2.2)=2.2^3-4.4-1=5.248>0$. The midpoint is $1.6$, and
\[
f(1.6)=1.6^3-3.2-1=4.096-4.2=-0.104<0,
\]
so the next interval is $[1.6,2.2]$. The root kept by the later cuts is the golden ratio $\varphi=\frac{1+\sqrt5}2$. It satisfies $\varphi^2=\varphi+1$, so $\varphi^3=\varphi\cdot\varphi^2=\varphi^2+\varphi=2\varphi+1$, and $\varphi^3-2\varphi-1=0$. You can see the identity without trusting the decimal $1.618$.

\begin{figure}[ht]
\centering
\fig{ch05-bisection}
\caption{Left: bisection for $f(x)=x^3-2x-1$ on $[1,2.2]$. Each step keeps the half on which $f$ changes sign. After six steps the root $\frac{1+\sqrt5}2$ sits in an interval shorter than $0.02$. Right: Rolle's theorem for $f(x)=x^3-3x$ on $[0,\sqrt3]$. The graph meets zero at both black points, and its lowest point on the segment is an interior rust point, where the tangent is horizontal.}
\label{fig:bisection}
\end{figure}

\begin{examplex}[three consequences]
Every positive real has a positive $n$th root, for each integer $n\ge1$. Every real polynomial of odd degree has a real root. Every continuous $f:[0,1]\to[0,1]$ has a fixed point.
\end{examplex}

For the root, $x\mapsto x^n-y$ is continuous, negative at $0$, and positive at $1+y$, because $(1+y)^n\ge1+ny>y$ by Bernoulli's inequality. The intermediate-value theorem supplies a positive root of $x^n=y$.

For the polynomial, take $P$ monic of odd degree $d$, and write
\[
P(x)=x^d\Bigl(1+\frac{a_{d-1}}x+\cdots+\frac{a_0}{x^d}\Bigr).
\]
Let $A=|a_{d-1}|+\cdots+|a_0|$. If $|x|>1$ and $|x|>2A$, each term $\bigl|\frac{a_j}{x^{d-j}}\bigr|$ is at most $\frac{|a_j|}{|x|}$, so the sum of those absolute values is at most $\frac A{|x|}<\frac12$. The parenthesis then exceeds $\frac12$. Since $d$ is odd, $P(x)$ has the sign of $x$ for all sufficiently large $|x|$. Theorem~\ref{thm:ivt} applies on a segment from a large negative number to a large positive one.

For the fixed point, set $g(x)=f(x)-x$. Then $g(0)=f(0)\ge0$ and $g(1)=f(1)-1\le0$. If either end is zero, that end is a fixed point. If $g(0)>0>g(1)$, then $-g(0)<0<-g(1)$, and Theorem~\ref{thm:ivt} gives a zero of $-g$, hence of $g$.

\begin{theorem}[extreme values]\label{thm:evt}
A continuous function on a segment $[a,b]$ is bounded, and it attains its supremum and its infimum.
\end{theorem}
\begin{proof}
Suppose $f$ were unbounded above. Then some points $x_n\in[a,b]$ would satisfy $f(x_n)\ge n$. Proposition~\ref{prop:seqcrit}(c) gives a subsequence $x_{n_k}\to d$. The limit $d$ lies in $[a,b]$, because inequalities survive limits and $a\le x_{n_k}\le b$. Continuity would give $f(x_{n_k})\to f(d)$, a finite number, while $f(x_{n_k})\ge n_k\to\infty$. That is impossible. So $f$ is bounded above, and $M=\sup_{[a,b]}f$ exists as a real number.

For each $n$, choose $x_n$ with $f(x_n)>M-\frac1n$. A subsequence converges to some $d\in[a,b]$, and
\[
M-\frac1{n_k}<f(x_{n_k})\le M
\]
forces $f(d)=M$ by the squeeze. Apply the same argument to $-f$ for the infimum.
\end{proof}

Both hypotheses earn their place. The function $\frac1x$ on $]0,1]$ is continuous and unbounded. The function $x$ on $]0,1[$ is bounded and attains neither bound. The function equal to $x$ on $[0,1[$ and equal to $0$ at $1$ is bounded on the segment $[0,1]$ and attains no maximum. Drop the segment, or drop continuity, and the theorem goes quiet.

\begin{think}
The function $x\mapsto x^2$ is continuous on $\R$, but a $\delta$ that works for a given $\varepsilon$ at the point $a$ has to shrink as $|a|$ grows, because the graph gets steeper. On the segment $[0,10]$, can one $\delta$ serve every point at once? What about $\frac1x$ on $]0,1]$?
\end{think}

\begin{theorem}[Heine: uniform continuity]\label{thm:heine}
A continuous function on a segment $[a,b]$ is \emph{uniformly continuous}: for every $\varepsilon>0$ there is a $\delta>0$ such that $|f(x)-f(y)|\le\varepsilon$ whenever $x,y\in[a,b]$ and $|x-y|\le\delta$.
\end{theorem}
\begin{proof}
Suppose not. Then some $\varepsilon>0$ defeats every $\delta$. Taking $\delta=\frac1n$, we find $x_n,y_n\in[a,b]$ with $|x_n-y_n|\le\frac1n$ and $|f(x_n)-f(y_n)|>\varepsilon$. Proposition~\ref{prop:seqcrit}(c) gives a subsequence $x_{n_k}\to c\in[a,b]$. Then $y_{n_k}\to c$ as well, because $|y_{n_k}-x_{n_k}|\le\frac1{n_k}\to0$. Continuity at $c$ gives
\[
f(x_{n_k})-f(y_{n_k})\to f(c)-f(c)=0,
\]
which contradicts $|f(x_{n_k})-f(y_{n_k})|>\varepsilon$.
\end{proof}

The difference is the order of the quantifiers. Continuity hands you a $\delta$ after you have named both $\varepsilon$ and the point. Uniform continuity hands you a $\delta$ after you have named only $\varepsilon$. On $]0,1]$ the function $\frac1x$ is continuous at each point and not uniformly continuous. The inputs $\frac1n$ and $\frac1{n+1}$ satisfy
\[
\Bigl|\frac1n-\frac1{n+1}\Bigr|=\frac1{n(n+1)}\to0,
\]
while their images stay $1$ apart. Figure~\ref{fig:uniform} is the pair $\frac13$ and $\frac12$: the inputs differ by $\frac16$, and the outputs differ by $1$. Later pairs are closer on the axis and still a unit apart on the vertical. Heine's theorem is what makes Riemann sums converge for every continuous integrand, in Chapter~\ref{ch:calculus}.

Eduard Heine published the theorem in 1872. The argument is already in Dirichlet's Berlin lectures of 1854 on definite integrals. Those lectures were printed only in 1904, which is why the theorem wears Heine's name.

\begin{figure}[ht]
\centering
\fig{ch05-uniform}
\caption{The graph of $\frac1x$. The inputs $\frac13$ and $\frac12$ differ by $\frac16$, drawn in ochre. Their images differ by $1$, drawn in rust. Pairs $\frac1n$ and $\frac1{n+1}$ get arbitrarily close as inputs and never get closer than $1$ as outputs, so no single $\delta$ serves $\varepsilon=\frac12$ on the whole of $]0,1]$.}
\label{fig:uniform}
\end{figure}

\begin{theorem}[Rolle, and the mean value theorem]\label{thm:mvt}
Let $f$ be continuous on $[a,b]$ and differentiable on $]a,b[$, with $a<b$.
\begin{enumerate}[label=(\alph*)]
\item \textbf{Fermat.} If $f$ has a maximum or a minimum on $]a,b[$ at the point $c$, then $f'(c)=0$.
\item \textbf{Rolle.} If $f(a)=f(b)$, then $f'(c)=0$ for some $c\in\,]a,b[$.
\item \textbf{Mean value.} $f(b)-f(a)=f'(c)\,(b-a)$ for some $c\in\,]a,b[$. Hence if $|f'|\le M$ on $]a,b[$, then $|f(x)-f(y)|\le M|x-y|$ on $[a,b]$. If $f'\ge0$ on $]a,b[$, then $f$ is non-decreasing on $[a,b]$. If $f'=0$ on $]a,b[$, then $f$ is constant on $[a,b]$.
\end{enumerate}
\end{theorem}
\begin{proof}
\emph{Fermat.} Suppose $c$ is a maximum. For $x>c$ close to $c$, the numerator $f(x)-f(c)$ is $\le0$ and the denominator $x-c$ is $>0$, so the difference quotient is $\le0$. Let $x\to c$ from the right: $f'(c)\le0$. For $x<c$ the numerator is still $\le0$ and the denominator is $<0$, so the quotient is $\ge0$. Let $x\to c$ from the left: $f'(c)\ge0$. Therefore $f'(c)=0$. A minimum of $f$ is a maximum of $-f$.

\emph{Rolle.} Theorem~\ref{thm:evt} says that $f$ attains its maximum and its minimum on $[a,b]$. If they are equal, $f$ is constant and every interior point works. If not, one of them differs from the common value $f(a)=f(b)$, so it is attained at an interior point, and Fermat applies there. The right half of Figure~\ref{fig:bisection} is this picture for $x^3-3x$ on $[0,\sqrt3]$: the values at the endpoints agree, and the interior minimum has a horizontal tangent.

\emph{Mean value.} Set
\[
g(x)=f(x)-f(a)-\frac{f(b)-f(a)}{b-a}(x-a).
\]
Then $g(a)=0$, and
\[
g(b)=f(b)-f(a)-\frac{f(b)-f(a)}{b-a}(b-a)=0.
\]
Rolle supplies a point $c$ with $g'(c)=0$, that is
\[
f'(c)=\frac{f(b)-f(a)}{b-a}.
\]
Apply the equality on any subsegment $[x,y]$ of $[a,b]$. If $|f'|\le M$, the absolute value of the difference of values is at most $M|y-x|$. If $f'\ge0$, the difference $f(y)-f(x)$ has the sign of $y-x$, so $f$ is non-decreasing. If $f'=0$, then $f(y)=f(x)$.
\end{proof}

Michel Rolle published the polynomial case in 1691, in a book on equations. The form for a general differentiable function uses Cauchy's definitions, which Rolle did not have. The mean value theorem is the sentence a speedometer has to satisfy. A car that covers $100$ kilometres in one hour has, at some instant, been travelling at exactly $100$ kilometres per hour. If the average speed on the journey exceeds a limit, the speed itself exceeds that limit somewhere on the road.

The next comparison is the one the polynomials chapter quotes when it estimates cotangents. The unit-circle facts used here are the ones the preface admits: $\sin0=\sin\pi=0$, $\sin\frac\pi2=1$, $\cos0=1$, $\cos\frac\pi2=0$, $\cos\pi=-1$, the number $\frac\pi2$ is the least positive zero of the cosine, and $\pi$ is the least positive angle at which the cosine equals $-1$. The addition law gives $\cos(x-x)=\cos^2x+\sin^2x=\cos0=1$.

\begin{proposition}[sine and tangent compared with the identity]\label{prop:sincmp}
On $[0,\pi]$, $\sin x\ge0$. For $0<x\le\frac\pi2$, one has $0<\sin x<x$. For $0<x<\frac\pi2$, one also has $x<\tan x$.
\end{proposition}
\begin{proof}
\emph{Sine increases on $[0,\frac\pi2]$.} The normalisation gives $\cos x>0$ on $\bigl[0,\frac\pi2\bigr)$. There $\sin'=\cos>0$ on the open interval. If $0\le x<y\le\frac\pi2$, the mean value theorem yields $\sin y-\sin x=\cos c\cdot(y-x)$ for some $c\in(x,y)\subset\bigl(0,\frac\pi2\bigr)$, so $\cos c>0$ and $\sin y>\sin x$. Sine is strictly increasing from $\sin0=0$ to $\sin\frac\pi2=1$. In particular $\sin x>0$ on $\bigl(0,\frac\pi2\bigr]$.

\emph{Sine stays nonnegative up to $\pi$.} Suppose some point of $\bigl(\frac\pi2,\pi\bigr)$ had a negative sine, and let $d$ be the infimum of such points. Then $d\in\bigl[\frac\pi2,\pi\bigr)$. The value $\sin d$ cannot be positive: continuity would keep sine positive on a neighbourhood of $d$, and $d$ would not be the infimum. It cannot be negative: continuity would keep sine negative on a neighbourhood of $d$, and some point of the set would lie to the left of $d$. So $\sin d=0$. Sine is already nonnegative on $\bigl[0,\frac\pi2\bigr]$, and it is nonnegative on $\bigl[\frac\pi2,d\bigr)$ because no point there belongs to the set whose infimum is $d$. Thus $\sin\ge0$ on $[0,d]$, so $\cos'=-\sin\le0$ on $(0,d)$. For $0\le x<y\le d$ the mean value theorem gives $\cos y-\cos x=-\sin c\cdot(y-x)\le0$. Cosine is non-increasing on $[0,d]$, and $\cos d\le\cos\frac\pi2=0$. Combined with $\sin d=0$ and $\sin^2d+\cos^2d=1$, this forces $\cos d=-1$. The leastness of $\pi$ then gives $d\ge\pi$, which contradicts $d<\pi$. So no such negative value exists, and $\sin x\ge0$ on $[0,\pi]$.

\emph{The two inequalities.} On $\bigl(0,\frac\pi2\bigr]$, sine is positive, so $\cos'=-\sin<0$ and cosine is strictly decreasing from $1$ to $0$. The mean value theorem gives $\sin x=\sin x-\sin0=x\cos c$ for some $c\in(0,x)$. Here $0<\cos c<1$, because $c\in\bigl(0,\frac\pi2\bigr)$. Thus $0<\sin x<x$. If also $x<\frac\pi2$, then $c<x<\frac\pi2$, so $\cos c>\cos x>0$. Therefore $\sin x=x\cos c>x\cos x$. Divide by $\cos x$: $x<\tan x$.
\end{proof}

Check $x=\frac12$. Then $\sin\frac12\approx0.479$ and $\tan\frac12\approx0.546$, and $0.479<\frac12<0.546$. The inequalities are not delicate at this particular point. The proof is there so that a later estimate may quote them.

Chapter~\ref{ch:calculus} builds the rest of the differential calculus on the mean value theorem: Taylor's formula with a remainder, and l'H\^opital's rule. Nothing here assumes a Taylor expansion.

\subsection*{Integrals, as far as the series need them}
The preface admits that a continuous real function on a closed bounded interval has an integral, with linearity, additivity over adjacent intervals, monotonicity, and $\int_a^b c\,\dd t=c(b-a)$. Reversed limits change the sign. Those properties compare sums with areas. The rule that computes the area is the next proposition. It uses continuity and the mean value theorem, and nothing else.

\begin{proposition}[a first fundamental theorem]\label{prop:early-ftc}
If $f$ is continuous on an interval $I$ and $a\in I$, then $F(x)=\int_a^x f(t)\,\dd t$ has derivative $F'(x)=f(x)$ at every interior point of $I$. Consequently, if $G$ is continuously differentiable on $[a,b]$, then $\int_a^b G'(t)\,\dd t=G(b)-G(a)$.
\end{proposition}
\begin{proof}
Let $x$ be interior, and let $h\ne0$ be small enough that the segment from $x$ to $x+h$ lies in $I$. Additivity gives
\[
\frac{F(x+h)-F(x)}h-f(x)=\frac1h\int_x^{x+h}\bigl(f(t)-f(x)\bigr)\,\dd t.
\]
Given $\varepsilon>0$, continuity at $x$ provides $\delta>0$ such that $|t-x|\le\delta$ implies $|f(t)-f(x)|\le\varepsilon$. Take $0<|h|\le\delta$.

If $h>0$, then $-\varepsilon\le f(t)-f(x)\le\varepsilon$ throughout $[x,x+h]$, so monotonicity of the integral yields $-\varepsilon h\le\int_x^{x+h}(f-f(x))\le\varepsilon h$. Divide by $h$: the difference quotient for $F$ lies within $\varepsilon$ of $f(x)$.

If $h=-k$ with $k>0$, the reversed-limit rule gives
\[
\frac1h\int_x^{x+h}(f-f(x))=\frac1k\int_{x-k}^x(f-f(x)),
\]
and the same bound on $[x-k,x]$ shows that this average lies within $\varepsilon$ of $0$. So $F'(x)=f(x)$.

For the consequence, the function $H(x)=\int_a^x G'(t)\,\dd t-G(x)$ has derivative $G'(x)-G'(x)=0$. Part (c) of Theorem~\ref{thm:mvt} says that $H$ is constant. At $x=a$ the integral vanishes, so the constant is $-G(a)$. Therefore $\int_a^x G'=G(x)-G(a)$. Set $x=b$.
\end{proof}

\section{Series}

A series is not a new kind of object. It is the sequence of its partial sums $S_n=u_0+\cdots+u_n$, and it \emph{converges} when $(S_n)$ does. The sum means that limit. Every test from Proposition~\ref{prop:seqcrit} applies to $(S_n)$. The new danger is rearrangement.

\begin{proposition}[series of non-negative terms]\label{prop:posseries}
Let $u_n,v_n\ge0$.
\begin{enumerate}[label=(\alph*)]
\item $\sum u_n$ converges if and only if its partial sums are bounded.
\item \textbf{Comparison.} If $u_n\le v_n$ and $\sum v_n$ converges, so does $\sum u_n$. If $u_n\sim v_n$, the two series converge or diverge together.
\item \textbf{Integral test.} If $f\ge0$ is continuous and non-increasing on $[1,\infty[$, then $\sum f(n)$ converges if and only if $\int_1^N f$ stays bounded as $N\to\infty$.
\item \textbf{Riemann.} $\sum\frac1{n^\alpha}$ converges if and only if $\alpha>1$.
\item \textbf{Ratio test.} If $u_n>0$ and $\frac{u_{n+1}}{u_n}\to\ell$, the series converges if $\ell<1$ and diverges if $\ell>1$.
\end{enumerate}
\end{proposition}
\begin{proof}
\emph{Partial sums.} They are non-decreasing. Proposition~\ref{prop:seqcrit}(a) is the claim.

\emph{Comparison.} The partial sums of $\sum u_n$ are at most those of $\sum v_n$, up to a shift of finitely many terms if the inequality $u_n\le v_n$ starts late. Finitely many terms never decide convergence. If $u_n\sim v_n$, then $\bigl|\frac{u_n}{v_n}-1\bigr|\le\frac12$ for $n$ large, so $\frac12 v_n\le u_n\le\frac32 v_n$. Compare in both directions.

\emph{Integral test.} On $[k,k+1]$ the function sits between its values at the ends: $f(k+1)\le f(t)\le f(k)$. Integrate and use $\int_k^{k+1}1\,\dd t=1$:
\[
f(k+1)\le\int_k^{k+1}f\le f(k).
\]
Figure~\ref{fig:integral} draws the comparison. Sum from $k=1$ to $N-1$:
\[
\sum_{k=2}^{N}f(k)\le\int_1^N f\le\sum_{k=1}^{N-1}f(k).
\]
The partial sums of the series and the integrals are bounded together.

\emph{Riemann's series.} If $\alpha\le0$, then $\frac1{n^\alpha}\not\to0$, so the series diverges. If $\alpha>0$, apply the integral test to $t^{-\alpha}$. The antiderivative is $\ln t$ when $\alpha=1$, and $\frac{t^{1-\alpha}}{1-\alpha}$ otherwise, by Proposition~\ref{prop:early-ftc} and the power rule already proved. From $1$ to $N$ this is $\ln N$, or $\frac{N^{1-\alpha}-1}{1-\alpha}$. The logarithm tends to infinity. If $\alpha<1$, then $1-\alpha>0$ and $N^{1-\alpha}$ tends to infinity. If $\alpha>1$, then $N^{1-\alpha}$ tends to $0$ and the integral tends to $\frac1{\alpha-1}$. Boundedness holds exactly for $\alpha>1$.

\emph{Ratio test.} If $\ell<1$, choose $r$ with $\ell<r<1$. For $n\ge N$, $\frac{u_{n+1}}{u_n}\le r$, so $u_n\le u_N r^{n-N}$. The comparison series is geometric with ratio $r$. If $\ell>1$, choose $r$ with $1<r<\ell$. Eventually $\frac{u_{n+1}}{u_n}\ge r$, so the terms are eventually increasing and bounded below by the positive number $u_N$. They do not tend to $0$, and the partial sums are unbounded.
\end{proof}

When the ratio tends to $1$, the test is silent. Both $\sum\frac1n$ and $\sum\frac1{n^2}$ have ratios tending to $1$, and they disagree about convergence. The direct comparison with a Riemann series is the repair: if $n^\alpha u_n$ stays bounded for some $\alpha>1$, the series converges; if $nu_n\ge c>0$ for $n$ large, it diverges. Exercise~\ref{ex:condensation} is Cauchy's condensation test, which reaches the series that live between $\frac1n$ and $\frac1{n^{1+\delta}}$.

\begin{figure}[ht]
\centering
\fig{ch05-integral}
\caption{A decreasing positive function. On each unit interval the graph sits between the rust left bar, of height $f(k)$, and the ochre right bar, of height $f(k+1)$. The integral over that interval is caught between the two areas, so the series $\sum f(n)$ and the integrals $\int_1^N f$ grow up together or stay bounded together.}
\label{fig:integral}
\end{figure}

\begin{examplex}[three series, three tests]
The series $\sum\frac{n^3}{3^n}$ converges, $\sum\frac{\ln n}{n^2}$ converges, and $\sum\sin\frac1n$ diverges.
\end{examplex}

The ratio of successive terms in the first series is $\frac{(n+1)^3}{3n^3}\to\frac13<1$. The ratio test applies.

For the logarithm, the mean value theorem on $[1,u]$ gives $\ln u-\ln1=\frac1c(u-1)$ for some $c\in(1,u)$, if $u>1$. Since $c>1$, $\ln u\le u-1\le u$. Thus $\ln n=4\ln(n^{1/4})\le4n^{1/4}$ for $n\ge1$, the case $n=1$ being $0=0$. Therefore
\[
0\le\frac{\ln n}{n^2}\le\frac{4}{n^{7/4}}.
\]
Here $\frac74>1$, so the comparison series converges, and so does $\sum\frac{\ln n}{n^2}$. At $n=16$, $\ln16\approx2.77$ and $4\cdot16^{1/4}=8$, so the inequality has room.

For the sine, $\sin\frac1n\sim\frac1n$ as $n\to\infty$, by the reminder on equivalents. The harmonic series diverges, so this one does too. The same comparison shows that $\sum\frac1{\sqrt{n(n+1)}}$ diverges, even though its general term is smaller than $\frac1n$: the ratio of the two terms tends to $1$.

\begin{think}
In a finite sum the order of the terms does not matter. Should $1-\frac12+\frac13-\frac14+\cdots$ keep that privilege? What could go wrong if the positive terms alone added to $+\infty$ and the negative terms alone to $-\infty$?
\end{think}

\begin{proposition}[signed series]\label{prop:signedseries}
\begin{enumerate}[label=(\alph*)]
\item If $\sum|u_n|$ converges, then $\sum u_n$ converges.
\item \textbf{Alternating series.} If $(a_n)$ is non-increasing and tends to $0$, then $\sum(-1)^na_n$ converges. The remainder $R_n=\sum_{k>n}(-1)^ka_k$ lies between $0$ and the first omitted term, and $|R_n|\le a_{n+1}$.
\item \textbf{Cauchy product.} If $\sum a_n$ and $\sum b_n$ converge absolutely, and $c_n=\sum_{k=0}^n a_kb_{n-k}$, then $\sum c_n$ converges absolutely and equals $\bigl(\sum a_n\bigr)\bigl(\sum b_n\bigr)$.
\end{enumerate}
\end{proposition}
\begin{proof}
\emph{Absolute convergence.} Let $S_n$ be the partial sums of $\sum u_n$. For $p>q$,
\[
|S_p-S_q|\le\sum_{q<k\le p}|u_k|.
\]
The right side is small for large $q$, because the partial sums of $\sum|u_k|$ form a convergent, hence Cauchy, sequence. Proposition~\ref{prop:seqcrit}(d) says that $(S_n)$ converges.

\emph{Alternating series.} First, $a_n\ge0$. If some $a_N<0$, then $a_n\le a_N<0$ for every $n\ge N$, and $(a_n)$ could not tend to $0$. Let $S_n=\sum_{k\le n}(-1)^ka_k$. Then
\begin{align*}
S_{2n+2}-S_{2n}&=-a_{2n+1}+a_{2n+2}\le0,\\
S_{2n+3}-S_{2n+1}&=a_{2n+2}-a_{2n+3}\ge0,
\end{align*}
because $(a_n)$ is non-increasing. So $(S_{2n})$ is non-increasing and $(S_{2n+1})$ is non-decreasing. Also $S_{2n}-S_{2n+1}=a_{2n+1}\to0$, and $S_{2n+1}\le S_{2n}$. The two subsequences are adjacent, so they share a limit $S$, and $S_{2n+1}\le S\le S_{2n}$. Every index is even or odd, so $S_n\to S$.

The limit sits between any two consecutive partial sums. Therefore
\[
|R_n|=|S-S_n|\le|S_{n+1}-S_n|=a_{n+1},
\]
and $R_n=S-S_n$ lies between $0$ and $S_{n+1}-S_n=(-1)^{n+1}a_{n+1}$.

\emph{Cauchy product.} Let $A_N=\sum_{n\le N}|a_n|$ and $B_N=\sum_{n\le N}|b_n|$, and write $A=\sum|a_n|$, $B=\sum|b_n|$. The triangle $k+l\le N$ sits inside the square $0\le k,l\le N$, which sits inside the triangle $k+l\le2N$. Figure~\ref{fig:cauchyprod} is the case $N=6$. Hence
\[
\sum_{n\le N}|c_n|\le\sum_{k+l\le N}|a_k||b_l|\le A_NB_N.
\]
The partial sums of $\sum|c_n|$ are bounded, so that series converges, and $\sum c_n$ converges absolutely by part (a).

For the value of the sum, compare the square with the large triangle:
\[
\Bigl|\sum_{n\le2N}c_n-\Bigl(\sum_{k\le N}a_k\Bigr)\Bigl(\sum_{l\le N}b_l\Bigr)\Bigr|
\le\sum_{\substack{k+l\le2N\\ \max(k,l)>N}}|a_k||b_l|
\le\Bigl(\sum_{k>N}|a_k|\Bigr)B+A\Bigl(\sum_{l>N}|b_l|\Bigr).
\]
The right side tends to $0$ as $N\to\infty$. The even partial sums of $\sum c_n$ therefore tend to $\bigl(\sum a_k\bigr)\bigl(\sum b_l\bigr)$. A convergent series has only one limit, so the full sequence of partial sums tends to the same product.
\end{proof}

The first even partial sums of $1-\frac12+\frac13-\frac14+\cdots$ show the alternating pattern before the notation does. With $a_n=\frac1{n+1}$,
\[
S_0=1,\quad S_2=\frac56,\quad S_4=\frac{47}{60},\quad S_1=\frac12,\quad S_3=\frac7{12},\quad S_5=\frac{37}{60}.
\]
The even sums $1$, $0.833$, $0.783$ fall. The odd sums $0.5$, $0.583$, $0.617$ rise. The gap between $S_4$ and $S_5$ is $\frac1{6}$, the first omitted term after $S_4$. They are heading for the same place. Exercise~\ref{ex:ln2} names it: $\ln2$.

The complex forms of these tests are in Chapter~\ref{ch:complex}, once a complex sequence has a real and an imaginary part. The proofs there quote the real case.

\begin{figure}[ht]
\centering
\fig{ch05-cauchyprod}
\caption{The Cauchy product for $N=6$. Each dot $(k,l)$ carries $a_kb_l$. The partial sum $\sum_{n\le N}c_n$ adds the charcoal triangle $k+l\le N$. The product of the two partial sums adds the grey square. The rust dots lie in the large triangle $k+l\le2N$ and outside the square, so each of them has $k>N$ or $l>N$. Absolute convergence makes their total small.}
\label{fig:cauchyprod}
\end{figure}

\begin{proposition}[when changing the order is safe]\label{prop:safe-sums}
For $a_{ij}\ge0$, the iterated sums $\sum_i\sum_j a_{ij}$ and $\sum_j\sum_i a_{ij}$ are equal, allowing the value $+\infty$. An absolutely convergent real series keeps its sum under any reordering of its terms.
\end{proposition}
\begin{proof}
Let $M$ be the supremum of all finite rectangle sums $\sum_{i=0}^m\sum_{j=0}^n a_{ij}$, with the convention $M=+\infty$ if those sums are unbounded. For fixed $m$, increasing $n$ shows $\sum_{i=0}^m\sum_j a_{ij}\le M$. Increasing $m$ shows that the first iterated sum is at most $M$. Every finite rectangle sum is at most that iterated sum, because a rectangle sum is a partial sum of it. So the iterated sum equals $M$. The same argument with the indices swapped gives the other iterated sum.

For a rearrangement, let $S=\sum_{k\ge0}u_k$ and let $\sum|u_k|$ converge. Given $\varepsilon>0$, choose $K$ with $\sum_{k>K}|u_k|<\varepsilon$. Let $\sigma$ be a permutation of the indices, and choose $N$ so that the set $\{\sigma(0),\ldots,\sigma(N)\}$ contains $\{0,\ldots,K\}$. Any later rearranged partial sum contains those first $K+1$ original terms. Its difference from $S$ is a sum of original terms with indices greater than $K$, so the absolute value is at most $\varepsilon$. The rearranged partial sums tend to $S$.
\end{proof}

\portrait{riemann}{Bernhard Riemann (1826--1866). Photograph, 1863. Public domain, via Wikimedia Commons.}

Absolute convergence is the permission to treat an infinite sum as a finite one. Without it, the permission is revoked. The epigraph rearranges the alternating harmonic series and cuts the sum in half. Exercise~\ref{ex:rearr} finds the sum for every pattern that takes $p$ positive terms and then $q$ negative ones. Exercise~\ref{ex:riemannrearr} shows that if $\sum u_n$ converges and $\sum|u_n|$ diverges, then some rearrangement converges to any real number you name, and some rearrangement diverges to $+\infty$. The positive terms alone and the negative terms alone are both infinite piles. You can spend from either pile for as long as you like.

\begin{think}
The sum $H_n=1+\frac12+\cdots+\frac1n$ and the integral $\ln n=\int_1^n\frac{\mathrm dt}t$ look at the same function. Sketch the bars of height $\frac1k$ on $[k,k+1]$ against the curve $\frac1t$. Does $H_n-\ln n$ stay bounded? Does it rise, or fall?
\end{think}

\begin{proposition}[Euler's constant]\label{prop:gamma}
Let $H_n=\sum_{k=1}^n\frac1k$. There is a constant $\gamma$ such that $H_n=\ln n+\gamma+o(1)$. The decimal begins $0.57721$.
\end{proposition}
\begin{proof}
Set $u_n=H_n-\ln n$ and $v_n=H_n-\ln(n+1)$. On $[n,n+1]$ one has $\frac1{n+1}\le\frac1t\le\frac1n$. Integrate, using Proposition~\ref{prop:early-ftc} for the logarithm:
\[
\frac1{n+1}\le\ln\Bigl(1+\frac1n\Bigr)\le\frac1n.
\]
The left half says
\[
u_{n+1}-u_n=\frac1{n+1}-\ln\Bigl(1+\frac1n\Bigr)\le0,
\]
so $(u_n)$ is non-increasing. The same inequality at $n+1$ says $\ln\bigl(1+\frac1{n+1}\bigr)\le\frac1{n+1}$, so
\[
v_{n+1}-v_n=\frac1{n+1}-\ln\Bigl(1+\frac1{n+1}\Bigr)\ge0,
\]
and $(v_n)$ is non-decreasing. The gap is $u_n-v_n=\ln\bigl(1+\frac1n\bigr)\to0$, and $v_n\le u_n$. The sequences are adjacent, so they share a limit $\gamma$, and $v_n\le\gamma\le u_n$. In particular $u_n=\gamma+o(1)$, which is the claim.
\end{proof}

The first values already fall toward the limit and are not there yet. $H_4=\frac{25}{12}=2.0833$ and $\ln4\approx1.386$, so $u_4\approx0.697$. $H_8\approx2.718$ and $\ln8\approx2.079$, so $u_8\approx0.638$. The constant itself begins $0.57721$. Euler studied it in a paper on harmonic progressions written in 1731, read in 1734 and published in 1740. He printed six decimals, $0.577218$. The last digit is wrong: the expansion begins $0.5772156649$.

Geometrically, $v_n=H_n-\ln(n+1)$ is the area of the curved triangles between the bars and the graph of $\frac1t$, up to $n$. Figure~\ref{fig:gamma} shows those regions through $n=8$, with total area $H_8-\ln9\approx0.521$. As further bars are drawn, the area increases to $\gamma$. The right panel is a preview of Exercise~\ref{ex:rearr}: the same terms, three running sums.

\begin{figure}[ht]
\centering
\fig{ch05-gamma}
\caption{Left: the coral regions between the bars of height $\frac1k$ and the curve $\frac1t$ have total area $H_8-\ln9\approx0.521$. Further bars raise this area to $\gamma\approx0.57721$. Right: the first $60$ partial sums of the alternating harmonic series, in charcoal, and of two rearrangements from Exercise~\ref{ex:rearr}. One positive term and then two negative ones, in rust, settles lower. Two positive terms and then one negative, in ochre, settles higher. Same terms, three sums.}
\label{fig:gamma}
\end{figure}

\begin{theorem}[Stirling]\label{thm:stirling}
$n!\sim\sqrt{2\pi n}\,\Bigl(\dfrac ne\Bigr)^n$.
\end{theorem}
\begin{proof}
Exercise~\ref{ex:stirling} proves that $n!\sim C\sqrt n\,\bigl(\frac ne\bigr)^n$ for some constant $C>0$. Theorem~\ref{thm:stirlingC} in Chapter~\ref{ch:calculus} shows $C=\sqrt{2\pi}$, by Wallis's integrals. The constant is not identified in this chapter.
\end{proof}

The history is an exchange, and both names belong in it. Abraham de Moivre, a Huguenot who had settled in London, proved in 1730 that $n!\sim C\sqrt n\,(n/e)^n$ and computed $C$ from a slowly convergent series as about $2.5074$. James Stirling, in the \emph{Methodus Differentialis} of the same year, showed that the constant is $\sqrt{2\pi}$, whose decimal begins $2.506628$. In the supplement to the \emph{Miscellanea Analytica}, de Moivre quotes Stirling's series and credits him. The approximation has carried Stirling's name since. The numerical gap between $2.5074$ and $2.5066$ is small, and it is the whole of Stirling's contribution to the leading term. It is also the difference between a constant you can only approximate and a constant you can name.

\section{Summation by parts}

The alternating test wants the signs to flip at every term. The series $\sum\frac{\sin n}n$ and $\sum\frac{\cos(n\theta)}{\sqrt n}$ flip according to a rhythm that is not that simple. The tool is a discrete integration by parts. Abel used it in 1826, in his paper on the binomial series.

\portrait{abel}{Niels Henrik Abel (1802--1829). Portrait by Johan G\o rbitz, 1826. Public domain, via Wikimedia Commons.}

\begin{proposition}[Abel's summation by parts]\label{prop:abelsum}
Let $A_k=a_0+\cdots+a_k$, with $A_{-1}=0$. Then for $0\le m\le n$,
\[
\sum_{k=m}^n a_kb_k=A_nb_n-A_{m-1}b_m+\sum_{k=m}^{n-1}A_k(b_k-b_{k+1}).
\]
\end{proposition}
\begin{proof}
Write $a_k=A_k-A_{k-1}$. Then
\[
\sum_{k=m}^n a_kb_k=\sum_{k=m}^n A_kb_k-\sum_{k=m}^n A_{k-1}b_k.
\]
In the second sum, set $j=k-1$. It becomes $\sum_{j=m-1}^{n-1}A_jb_{j+1}$. Subtract. For each index $k$ with $m\le k\le n-1$, the coefficient of $A_k$ is $b_k-b_{k+1}$. The unpaired terms are $A_nb_n$ from the first sum and $-A_{m-1}b_m$ from the second.
\end{proof}

The identity moves the summing onto $(a_k)$, whose partial sums $A_k$ you hope are bounded, and leaves the differences $b_k-b_{k+1}$, which you hope are small. That is what $\int u'\,v=[uv]-\int u\,v'$ does with a derivative.

\begin{theorem}[Dirichlet's test]\label{thm:dirichlettest}
Let $(a_k)$ be a real sequence whose partial sums satisfy $|A_n|\le M$, and let $(b_k)$ be real, non-increasing, and convergent to $0$. Then $\sum a_kb_k$ converges, and
\[
\Bigl|\sum_{k>n}a_kb_k\Bigr|\le2Mb_{n+1}.
\]

\textbf{Abel's test.} If $\sum a_k$ converges and $(b_k)$ is monotone and bounded, then $\sum a_kb_k$ converges.
\end{theorem}
\begin{proof}
Since $(b_k)$ is non-increasing and tends to $0$, the same argument as in the alternating test gives $b_k\ge0$. Proposition~\ref{prop:abelsum}, for $m\le n$, yields
\begin{align*}
\Bigl|\sum_{k=m}^n a_kb_k\Bigr|
&\le Mb_n+Mb_m+M\sum_{k=m}^{n-1}(b_k-b_{k+1})\\
&=Mb_n+Mb_m+M(b_m-b_n)=2Mb_m.
\end{align*}
The differences $b_k-b_{k+1}$ are nonnegative, so the sum telescopes without absolute values. As $m\to\infty$, the bound $2Mb_m$ tends to $0$. The partial sums of $\sum a_kb_k$ are therefore Cauchy, and the series converges by Proposition~\ref{prop:seqcrit}(d).

Hold $m$ fixed and let $n\to\infty$. The block sum from $m$ to $n$ tends to the tail $\sum_{k\ge m}a_kb_k$, and the bound $2Mb_m$ does not depend on $n$. So the tail is at most $2Mb_m$ in absolute value. Set $m=n+1$.

For Abel's test, monotonicity and boundedness give $b_k\to b$ for some finite $b$. If $(b_k)$ is non-increasing, then $b_k-b$ is non-increasing and tends to $0$. If $(b_k)$ is non-decreasing, then $b-b_k$ is non-increasing and tends to $0$. Dirichlet's test applies to $a_k(b_k-b)$ or to $a_k(b-b_k)$. The series $\sum a_kb$ converges because $\sum a_k$ does. Add the two series.
\end{proof}

\begin{examplex}[the series $\sum\frac{\sin(k\theta)}k$]
For every real $\theta$, the series $\sum_{k\ge1}\frac{\sin(k\theta)}k$ converges.
\end{examplex}

If $\theta\in2\pi\Z$, every term is zero. Now suppose $\theta\notin2\pi\Z$, so $\sin\frac\theta2\ne0$. The product-to-sum identity
\[
2\sin\frac\theta2\sin(k\theta)=\cos\Bigl(\bigl(k-\tfrac12\bigr)\theta\Bigr)-\cos\Bigl(\bigl(k+\tfrac12\bigr)\theta\Bigr)
\]
comes from $2\sin A\sin B=\cos(A-B)-\cos(A+B)$, with $A=\frac\theta2$ and $B=k\theta$. Multiply the partial sum by $2\sin\frac\theta2$:
\[
2\sin\frac\theta2\sum_{k=1}^n\sin(k\theta)=\cos\frac\theta2-\cos\Bigl(\bigl(n+\tfrac12\bigr)\theta\Bigr),
\]
because the intermediate cosines cancel. The right side has absolute value at most $2$, so
\[
\Bigl|\sum_{k=1}^n\sin(k\theta)\Bigr|\le\frac1{\bigl|\sin\frac\theta2\bigr|}.
\]
The partial sums of $\sin(k\theta)$ are bounded. The sequence $\frac1k$ decreases to $0$. Dirichlet's test says that $\sum\frac{\sin(k\theta)}k$ converges.

The sum equals $\frac{\pi-\theta}2$ for $0<\theta<2\pi$. That identification is a Fourier series. Volume~5, \emph{Calculation: From Formula to Insight}, proves it. Exercise~\ref{ex:sinn} shows that at $\theta=1$ the convergence is not absolute. Riemann's warning applies: this series converges by cancellation.

\section{How many terms buy a decimal}

Convergence says that some tail is small. It does not say which tail. A decimal with no error bound is a guess that has put on a tie.

\begin{examplex}[a remainder for $\ln2$]
For every $n\ge1$,
\[
\ln2-\sum_{k=1}^n\frac{(-1)^{k-1}}k=(-1)^n\int_0^1\frac{t^n}{1+t}\,\dd t,
\]
and the integral lies between $\frac1{2(n+1)}$ and $\frac1{n+1}$.
\end{examplex}

Start from the finite geometric sum, not from an infinite series:
\[
\frac1{1+t}=\sum_{k=0}^{n-1}(-t)^k+\frac{(-t)^n}{1+t},\qquad t\ne-1.
\]
Integrate from $0$ to $1$. The left side is $\int_0^1\frac{\mathrm dt}{1+t}=\ln2$, by Proposition~\ref{prop:early-ftc}. The powers integrate to the alternating harmonic sum, because $\int_0^1 t^{k-1}\,\dd t=\frac1k$ after the sign $(-1)^{k-1}$ is restored, with the index shifted by one. The remainder integral is exactly the displayed formula. On $[0,1]$, $\frac12\le\frac1{1+t}\le1$, so
\[
\frac1{2(n+1)}\le\int_0^1\frac{t^n}{1+t}\,\dd t\le\int_0^1 t^n\,\dd t=\frac1{n+1}.
\]
The alternating-series bound $|R_n|\le\frac1{n+1}$ is honest. It is also slow. At $n=1000$ it promises an error of at most $\frac1{1001}$, which is not six correct decimals. Figure~\ref{fig:slow} shows the first twelve partial sums closing on $\ln2$ from both sides, at that same slow rate.

\begin{figure}[ht]
\centering
\fig{review05-slow}
\caption{The first twelve partial sums of the alternating harmonic series, approaching $\ln2$ from both sides. The dotted line is $\ln2$. The gap is of size about $\frac1n$. A thousand terms still leave an error near $10^{-3}$.}
\label{fig:slow}
\end{figure}

\subsection*{A faster route to the same number}
Six positive terms are enough for an error below $10^{-6}$, if the integral is chosen so that the powers decay.

The antiderivative of $\frac2{1-t^2}$ is $\ln\frac{1+t}{1-t}$, because
\[
\frac{\mathrm d}{\mathrm dt}\bigl(\ln(1+t)-\ln(1-t)\bigr)=\frac1{1+t}+\frac1{1-t}=\frac2{1-t^2}.
\]
At $t=\frac13$ this antiderivative equals $\ln\frac{4/3}{2/3}=\ln2$, and at $t=0$ it equals $0$. So
\[
\ln2=\int_0^{1/3}\frac2{1-t^2}\,\dd t.
\]
The upper limit is $\frac13$ because that is the point at which the antiderivative hits $\ln2$, not because of a later convenience.

Expand the integrand before integrating, and stop at a finite sum. For $|u|<1$,
\[
\frac1{1-u}=\sum_{k=0}^{m-1}u^k+\frac{u^m}{1-u}.
\]
Put $u=t^2$, multiply by $2$, and integrate from $0$ to $\frac13$:
\[
\ln2=2\sum_{k=0}^{m-1}\frac{3^{-(2k+1)}}{2k+1}+R_m,
\qquad
R_m=2\int_0^{1/3}\frac{t^{2m}}{1-t^2}\,\dd t.
\]
On $\bigl[0,\frac13\bigr]$ one has $t^2\le\frac19$, so $1-t^2\ge\frac89$ and $\frac1{1-t^2}\le\frac98$. Therefore
\[
0<R_m\le2\cdot\frac98\int_0^{1/3}t^{2m}\,\dd t=\frac9{4(2m+1)\,3^{2m+1}}.
\]
At $m=6$ the bound is $\frac1{9211644}$, which is less than $1.1\times10^{-7}$.

The six-term sum equals $\frac{15757912}{22733865}$. Its decimal begins $0.69314707$, so the fraction lies strictly between $0.69314707$ and $0.69314708$. The remainder is positive and smaller than $0.00000011$. Therefore
\[
0.69314707<\ln2<0.69314708+0.00000011=0.69314719<0.693148,
\]
and in particular $0.693147<\ln2<0.693148$. No infinite sum was integrated term by term. The finite identity carried an explicit remainder, and the remainder was estimated by a constant on a short interval. The alternating harmonic series was not short of patience. It was the wrong representation of the same number.

\section*{What this chapter prepares}
A sequence converges when every tail sits inside a band whose width was chosen in advance. Completeness turns that game into existence theorems: a bounded monotone sequence converges, a continuous function on a segment attains its bounds and its intermediate values, and a derivative that vanishes forces the function to be constant. A series converges when its partial sums do. Absolute convergence is the condition under which the sum survives a change of order. The comparison $0<\sin x<x<\tan x$ on $\bigl(0,\frac\pi2\bigr)$ is a theorem, available for later estimates.

The floor is now defined for every real number, which is what Dirichlet's approximation argument in Chapter~\ref{ch:counting} was waiting for. The hat-check proportions $D_n/n!$ are a sequence in the sense of this chapter. Whether they converge is a precise question. Naming the limit $1/e$ uses the exponential series, which is justified when the exponential is compared with its Taylor polynomial in Chapter~\ref{ch:calculus}.

The peaks in the proof of Bolzano--Weierstrass were a monotone piece of a bounded sequence, cut out by size. The next chapter asks what an order is. Among any $n^2+1$ distinct numbers, some $n+1$ of them, read from left to right, form a monotone run. Chapter~\ref{ch:counting} counted that fact with pigeonholes. Chapter~\ref{ch:orders} explains it as the theorem of Erd\H os and Szekeres on chains and antichains, and then asks when a process is forced to stop.

\section{Exercises}

\begin{exercise}[$\star$]\label{ex:nature}
Find the nature of $\sum\frac{n!}{n^n}$, and compute $\sum_{n\ge1}\frac{n^2}{2^n}$.
\end{exercise}

\begin{exercise}[$\star$]\label{ex:bertrand}
Show that $\sum_{n\ge2}\frac1{n\ln n}$ diverges and that $\sum_{n\ge2}\frac1{n\ln^2n}$ converges.
\end{exercise}

\begin{exercise}[$\star\star$]\label{ex:ln2}
Show that $\displaystyle\sum_{n=1}^\infty\frac{(-1)^{n+1}}n=\ln2$.
\end{exercise}

\begin{exercise}[$\star\star$]\label{ex:sin}
Let $u_0\in\left]0,\frac\pi2\right]$ and $u_{n+1}=\sin u_n$.
\begin{enumerate}[label=(\alph*)]
\item Show that $u_n\to0$.
\item \emph{Return after Taylor's theorem (Theorem~\ref{thm:taylor}).}
Show that $\frac1{u_{n+1}^2}-\frac1{u_n^2}\to\frac13$, and deduce that $u_n\sim\sqrt{\frac3n}$.
\end{enumerate}
\end{exercise}

\begin{exercise}[$\star\star$]\label{ex:kempner}
\emph{Kempner's series.} Show that the sum of $\frac1n$ over the integers $n\ge1$ whose decimal expansion contains no digit $9$ converges, and is less than $80$.
\end{exercise}

\begin{exercise}[$\star\star\star$]\label{ex:rearr}
\emph{Rearranging the alternating harmonic series.} Let $p,q\ge1$. Rearrange $\sum\frac{(-1)^{n+1}}n$ by taking, in turn, the next $p$ positive terms
$\frac11,\frac13,\frac15,\dots$, then the next $q$ negative terms $-\frac12,-\frac14,\dots$, and so on. Show that the new series converges to
\[
\ln2+\frac12\ln\frac pq.
\]
For $p=1$, $q=2$ this is the $\frac12\ln2$ of the opening.
\end{exercise}

\begin{exercise}[$\star\star\star$]\label{ex:stirling}
\emph{Stirling, up to a constant.} Let $d_n=\ln n!-\left(n+\frac12\right)\ln n+n$.
\begin{enumerate}[label=(\alph*)]
\item Show that $d_n-d_{n+1}=\left(n+\frac12\right)\ln\left(1+\frac1n\right)-1=\frac1{12n^2}+O\left(\frac1{n^3}\right)$.
\item Deduce that $(d_n)$ converges, and that $n!\sim C\sqrt n\left(\frac ne\right)^n$ for some $C>0$.
\end{enumerate}
\end{exercise}

\subsection*{Further exercises}

\begin{exercise}[$\star$]\label{ex:oresme}
\emph{Oresme, around 1350.}
\begin{enumerate}[label=(\alph*)]
\item Show that $H_{2^m}\ge1+\frac m2$ by grouping the terms $\frac1k$ for $2^{j-1}<k\le2^j$, and conclude again that $\sum\frac1n$ diverges, this time without integrals.
\item Show without integrals that $\sum\frac1{n^2}$ converges, with sum at most $2$: compare $\frac1{n^2}$ with $\frac1{n(n-1)}$ and telescope.
\item How many terms of the harmonic series are needed before the partial sum exceeds $10$? Estimate with Proposition~\ref{prop:gamma}.
\end{enumerate}
\end{exercise}

\begin{exercise}[$\star$]\label{ex:nested}
\emph{Nested radicals.} Let $c>0$, $u_0=0$ and $u_{n+1}=\sqrt{c+u_n}$. Show that $(u_n)$ converges and find its limit. What is $\sqrt{2+\sqrt{2+\sqrt{2+\cdots}}}$? Show that for $c=2$, $u_n=2\cos\frac\pi{2^{n+1}}$, and deduce Vi\`ete's formula of 1593,
\[
\frac2\pi=\frac{\sqrt2}2\cdot\frac{\sqrt{2+\sqrt2}}2\cdot\frac{\sqrt{2+\sqrt{2+\sqrt2}}}2\cdots
\]
\emph{(For the last step, use $\sin x=2^n\sin\frac x{2^n}\prod_{k=1}^n\cos\frac x{2^k}$ and $\frac{\sin t}t\to1$.)}
\end{exercise}

\begin{exercise}[$\star\star$]\label{ex:newton}
\emph{Newton's method for $\sqrt2$.} Let $x_0=1$ and $x_{n+1}=\frac12\left(x_n+\frac2{x_n}\right)$.
\begin{enumerate}[label=(\alph*)]
\item Show that $x_{n+1}-\sqrt2=\frac{(x_n-\sqrt2)^2}{2x_n}$, hence $0\le x_{n+1}-\sqrt2\le\frac{(x_n-\sqrt2)^2}{2\sqrt2}$ for $n\ge1$.
\item Deduce that $0\le x_n-\sqrt2\le2\sqrt2\left(\frac{x_1-\sqrt2}{2\sqrt2}\right)^{2^{n-1}}$ for $n\ge1$: the number of correct decimals roughly doubles at each step. How many steps give $100$ decimals?
\item Explain the recipe: $x_{n+1}$ is the point where the tangent to $y=x^2-2$ at $x_n$ meets the axis.
\end{enumerate}
\end{exercise}

\begin{exercise}[$\star\star$]\label{ex:stolz}
\emph{Stolz--Ces\`aro.} Let $(b_n)$ be strictly increasing with $b_n\to+\infty$, and suppose $\frac{a_{n+1}-a_n}{b_{n+1}-b_n}\to\ell\in\R$. Show that $\frac{a_n}{b_n}\to\ell$.
Deduce that $\frac{1^k+2^k+\dots+n^k}{n^{k+1}}\to\frac1{k+1}$ for every integer $k\ge0$, and that $\frac{\ln n!}{n\ln n}\to1$.
\end{exercise}

\begin{exercise}[$\star\star$]\label{ex:condensation}
\emph{Cauchy's condensation test.} Let $(u_n)$ be non-increasing with $u_n\ge0$. Show that $\sum u_n$ converges if and only if $\sum_k2^ku_{2^k}$ does. Recover
the nature of $\sum\frac1{n^\alpha}$, and show that $\sum_{n\ge3}\frac1{n\ln n\,(\ln\ln n)^\beta}$ converges if and only if $\beta>1$.
\end{exercise}

\begin{exercise}[$\star\star$]\label{ex:sinn}
Show that $\sum_{n\ge1}\frac{\sin n}n$ converges but $\sum_{n\ge1}\frac{|\sin n|}n$ diverges. \emph{(Use $|\sin n|\ge\sin^2n=\frac{1-\cos2n}2$.)}
\end{exercise}

\begin{exercise}[$\star\star\star$]\label{ex:riemannrearr}
\emph{Riemann's rearrangement theorem (1854).} Let $\sum u_n$ be a real series that converges but not absolutely.
\begin{enumerate}[label=(\alph*)]
\item Let $p_n=\max(u_n,0)$ and $q_n=\max(-u_n,0)$, so that $u_n=p_n-q_n$ and $|u_n|=p_n+q_n$. Show that $\sum p_n$ and $\sum q_n$ both diverge to $+\infty$.
\item Let $S\in\R$. Build a rearrangement of the $u_n$ that converges to $S$: take positive terms, in order, until the partial sum first exceeds $S$; then negative terms until it first drops below $S$; then positive terms again, and so on. Insert the zero terms, in their original order, one after each phase while any remain. Show that every term is used exactly once and that the partial sums converge to $S$.
\item Show that there are also rearrangements that diverge to $+\infty$.
\end{enumerate}
\end{exercise}

\begin{exercise}[$\star\star\star$]\label{ex:eulerbasel}
\emph{Euler's shortcut to $\zeta(2)$ (1731).} Let $\zeta(2)=\sum_{n\ge1}\frac1{n^2}$; this exercise does not use the value $\frac{\pi^2}6$. For $0\le x<1$ let
$L(x)=\displaystyle\int_0^x\frac{-\ln(1-t)}t\dd t$ (the integrand extends continuously by $1$ at $t=0$).
\begin{enumerate}[label=(\alph*)]
\item From $\frac1{1-u}=\sum_{k=0}^{n-1}u^k+\frac{u^n}{1-u}$, show that $L(x)=\sum_{k=1}^\infty\frac{x^k}{k^2}$ for $0\le x<1$, and that $L(x)\to\zeta(2)$ as $x\to1^-$.
\item Show that $L(x)+L(1-x)+\ln x\ln(1-x)$ is constant on $]0,1[$, equal to $\zeta(2)$.
\item Deduce Euler's formula $\zeta(2)=(\ln2)^2+2\sum_{n\ge1}\frac1{n^22^n}$, and show that twenty terms of the last series give $\zeta(2)$ to eight decimals, whereas twenty terms of $\sum\frac1{n^2}$ give one.
\end{enumerate}
\end{exercise}

\begin{challenge}
\begin{enumerate}[label=(\arabic*)]
\item \emph{Is $\gamma$ irrational?} Nobody knows. Assume a certified computation has established at least $1000$ positive partial quotients after the integer part of $\gamma$. Show that if $\gamma=p/q$ in lowest terms then $q\ge F_{1001}>10^{208}$. Use the recurrence $q_j=a_jq_{j-1}+q_{j-2}$ for convergent denominators, with $q_{-1}=0$, $q_0=1$ and $a_j\ge1$. Much stronger published bounds use the actual partial quotients, not merely their number.
\item \emph{The Erd\H{o}s--Borwein constant.} Show that $E=\sum_{n\ge1}\frac1{2^n-1}$ converges and equals $\sum_{n\ge1}\frac{d(n)}{2^n}$, where $d(n)$ is the number of divisors of $n$. Erd\H{o}s proved in 1948 that $E$ is irrational; his proof is three pages of elementary mathematics, and a good first research paper to read.
\item \emph{Kempner with other digits.} Show that for every digit $d$ the series of $\frac1n$ over the integers whose decimal expansion avoids $d$ converges, and that the same holds for integers avoiding any fixed block of digits, such as $42$. Why does the sum over integers avoiding the block $42$ (about $228.4$) dwarf the sum over integers avoiding the digit $9$ (about $22.9$)?
\end{enumerate}
Notes at the end of the chapter.
\end{challenge}

\begin{beauty}
Exercise~\ref{ex:riemannrearr} is Riemann's theorem. A series that converges, but not absolutely, can be rearranged so that it converges to any real number you name. Absolute convergence is the condition under which an infinite sum keeps the habits of a finite one. The same condition returns when an expectation is an infinite sum, in Volume~4, \emph{Chance}, and when integrals are allowed to change order, in Volume~7, \emph{Spaces}. The five theorems on a segment have a similar moral. Over $\Q$, a continuous function can change sign and never hit zero. One axiom, the existence of a least upper bound, is what makes the picture true on $\R$. Euler's constant $\gamma$ is still unsettled in a narrower way: nobody knows whether it is irrational. Challenge~(1) turns a long string of continued-fraction digits into a lower bound on the denominator of any rational that $\gamma$ might be.
\end{beauty}

\begin{exercise}[$\star$]\label{pb:schwartz}
\emph{Schwartz--Zippel.} Let $F$ be a field and $S\subset F$ finite.
\begin{enumerate}[label=(\alph*)]
\item Let $f\in F[x_1,\dots,x_n]$ be non-zero of total degree $d$, and $r$ uniform in $S^n$. Show that $P(f(r)=0)\le\frac d{|S|}$,
and that the bound is reached.
\item (Edmonds, 1967.) Let $G$ be a bipartite graph on $\{u_1,\dots,u_n\}\cup\{v_1,\dots,v_n\}$, and $X$ the $n\times n$ matrix with
$X_{ij}=x_{ij}$ (a variable) if $u_iv_j$ is an edge and $0$ otherwise. Show that $\det X\ne0$ in $\Q[x_{ij}]$ if and only if $G$ has
a perfect matching.
\item Deduce a randomised test for perfect matchings that uses one determinant over $\mathbb F_p$ and errs with probability at
most $\frac np$, and only in one direction. Which one?
\end{enumerate}
\end{exercise}

\begin{exercise}[$\star\star$]\label{pb:nullstellensatz}
\emph{Alon's Combinatorial Nullstellensatz.} Let $F$ be a field and $f\in F[x_1,\dots,x_n]$.
\begin{enumerate}[label=(\alph*)]
\item Let $S_1,\dots,S_n\subset F$ be finite. Show that if $\deg_{x_i}f<|S_i|$ for every $i$ and $f$ vanishes on
$S_1\times\dots\times S_n$, then $f=0$.
\item (Alon, 1999.) Suppose that $\deg f=t_1+\dots+t_n$ and that the coefficient of $x_1^{t_1}\cdots x_n^{t_n}$ in $f$ is non-zero.
Show that if $|S_i|>t_i$ for every $i$, then $f(s)\ne0$ for some $s\in S_1\times\dots\times S_n$.
\item (Alon's permanent lemma.) Let $A\in M_n(F)$ have non-zero permanent $\operatorname{per}A=\sum_\sigma\prod_iA_{i\sigma(i)}$,
let $b\in F^n$, and let $S_1,\dots,S_n\subset F$ have two elements each. Show that there is $x\in S_1\times\dots\times S_n$ such
that $(Ax)_i\ne b_i$ for every $i$.
\end{enumerate}
\end{exercise}

\begin{exercise}[$\star\star$]\label{pb:newton}
\emph{Newton's inequalities.} Let $P(x)=\sum_{k=0}^na_kx^k$ be a real polynomial of degree $n$ all of whose roots are real.
\begin{enumerate}[label=(\alph*)]
\item Show that for $1\le k\le n-1$,
\[
a_k^2\ge a_{k-1}a_{k+1}\Bigl(1+\frac1k\Bigr)\Bigl(1+\frac1{n-k}\Bigr).
\]
\item Deduce that the numbers $c(n,k)$ of permutations of $[n]$ with $k$ cycles form a log-concave, hence unimodal, sequence in $k$.
\end{enumerate}
\end{exercise}

\section{Hints}

\begin{hint}{ex:nature}
Ratio test for the first series. For the second, call the sum $B$ and write $n=m+1$: you get $B$ back, plus a multiple of the easier sum $\sum\frac n{2^n}$, which yields to the same trick.
\end{hint}
\begin{hint}{ex:bertrand}
Compare with integrals: what are the derivatives of $\ln\ln t$ and of $\frac1{\ln t}$?
\end{hint}
\begin{hint}{ex:ln2}
$S_{2n}=H_{2n}-2\cdot\frac12H_n$: all terms minus twice the even ones. Then use $H_n=\ln n+\gamma+o(1)$.
\end{hint}
\begin{hint}{ex:sin}
(a) $\sin x<x$ for $x>0$. (b) After Theorem~\ref{thm:taylor}, expand $\frac1{\sin^2x}-\frac1{x^2}$ as $x\to0$, then use the corollary of Ces\`aro.
\end{hint}
\begin{hint}{ex:kempner}
Count the admissible numbers with $k$ digits, and bound each of their reciprocals by $10^{-(k-1)}$.
\end{hint}
\begin{hint}{ex:rearr}
After $n$ blocks the partial sum is $\sum_{j\le pn}\frac1{2j-1}-\sum_{j\le qn}\frac1{2j}$. Write both with harmonic numbers. Then check the partial sums inside a block.
\end{hint}
\begin{hint}{ex:stirling}
(a) $d_n-d_{n+1}$ simplifies because $\ln(n+1)!-\ln n!=\ln(n+1)$.
Integrate $1/(1+t)=1-t+t^2-t^3/(1+t)$ from $0$ to $1/n$ and bound the remainder. (b) A series whose terms are $O\left(\frac1{n^2}\right)$ converges absolutely.
\end{hint}
\begin{hint}{ex:oresme}
(a) The block $2^{j-1}<k\le2^j$ has $2^{j-1}$ terms, each at least $\frac1{2^j}$. (b) $\frac1{n^2}<\frac1{n-1}-\frac1n$ for $n\ge2$. (c) Solve $\ln n+\gamma\approx10$.
\end{hint}
\begin{hint}{ex:nested}
The limit $\ell$ satisfies $\ell=\sqrt{c+\ell}$. For monotonicity, show by induction that $u_n<u_{n+1}<\ell$. For $c=2$, use $2+2\cos x=4\cos^2\frac x2$.
\end{hint}
\begin{hint}{ex:newton}
(a) Compute $x_{n+1}-\sqrt2$ with a common denominator and recognise a square. (b) Put $e_n=\frac{x_n-\sqrt2}{2\sqrt2}$ and show $e_{n+1}\le e_n^2$.
\end{hint}
\begin{hint}{ex:stolz}
For $n>N$ write $a_n-a_N$ as a telescoping sum of the differences $a_{k+1}-a_k$, each close to $\ell(b_{k+1}-b_k)$; then divide by $b_n$ and let $n\to\infty$ with $N$ fixed, as in the proof of Ces\`aro.
\end{hint}
\begin{hint}{ex:condensation}
Group the terms $u_n$ for $2^k\le n<2^{k+1}$: there are $2^k$ of them, between $u_{2^{k+1}}$ and $u_{2^k}$. Apply the test twice for the last series.
\end{hint}
\begin{hint}{ex:sinn}
Convergence: the example of the section with $\theta=1$. Divergence: $\sum\frac{\cos2n}n$ converges by Dirichlet's test, and $\sum\frac1{2n}$ does not.
\end{hint}
\begin{hint}{ex:riemannrearr}
(a) If $\sum p_n$ converged, so would $\sum q_n=\sum(p_n-u_n)$, and then $\sum|u_n|$. (b) Each phase ends because the positive (or negative) terms alone have infinite sum; after each change of direction the partial sum is within one term of $S$, and the terms tend to $0$. (c) Aim at $1$, then $2$, then $3$, and so on.
\end{hint}
\begin{hint}{ex:eulerbasel}
(a) Integrate the identity in $u$ from $0$ to $t$, divide by $t$, integrate from $0$ to $x$; bound the remainder using $\frac1{1-u}\le\frac1{1-x}$ for $u\le t\le x$; the limit $x\to1$ is a sandwich between partial sums. (b) Differentiate, with the fundamental theorem of calculus (Proposition~\ref{prop:early-ftc}); then let $x\to0$. (c) Put $x=\frac12$; the tail of $\sum\frac1{n^22^n}$ after $20$ terms is at most $\frac1{21^2}\cdot\frac1{2^{20}}$.
\end{hint}

\begin{hint}{pb:schwartz}
(a) Induction on $n$. Write $f$ as a polynomial in $x_1$ whose coefficients are polynomials in the other variables:
\[
f=\sum_{i\le k}x_1^i\,g_i(x_2,\dots,x_n),\qquad g_k\ne0.
\]
\end{hint}

\begin{hint}{pb:nullstellensatz}
(b) Shrink the $S_i$ to size $t_i+1$, and divide $f$ by the polynomials $\prod_{s\in S_i}(x_i-s)$. If $f$ vanished on the grid, the
remainder would vanish by (a).
\end{hint}

\begin{hint}{pb:newton}
Derivatives and reversals of polynomials with real roots have real roots. Reduce to a quadratic.
\end{hint}

\section{Solutions}
\begin{solution}{ex:nature}
\idea Factorials and powers $n^n$ are made for the ratio test, because the ratio $\frac{u_{n+1}}{u_n}$ simplifies dramatically. For $\sum\frac{n^2}{2^n}$, once convergence is known (ratio test again), the sum can be found without any calculus: shifting the index by one turns $n^2$ into $(n+1)^2=n^2+2n+1$, which brings back the sum itself plus simpler sums. So first compute the simpler sum $\sum\frac n{2^n}$ by the same trick.
\writeup With $u_n=\frac{n!}{n^n}$, $\frac{u_{n+1}}{u_n}=\frac{(n+1)\,n^n}{(n+1)^{n+1}}=\left(1+\frac1n\right)^{-n}\to\frac1e<1$: the series converges by the ratio test.
The series $A=\sum_{n\ge1}\frac n{2^n}$ and $B=\sum_{n\ge1}\frac{n^2}{2^n}$ converge by the ratio test (the ratios tend to $\frac12$). Shifting the index ($n=m+1$),
\[A=\sum_{m\ge0}\frac{m+1}{2^{m+1}}=\frac12\Big(A+\sum_{m\ge0}\frac1{2^m}\Big)=\frac12(A+2),\]
so $A=2$. In the same way, with $(m+1)^2=m^2+2m+1$,
\[B=\sum_{m\ge0}\frac{(m+1)^2}{2^{m+1}}=\frac12\,(B+2A+2),\]
so $B=2A+2=6$. (If you remember that power series can be differentiated termwise inside their disc of convergence, a fact this volume does not prove, you could instead apply $x\frac{\mathrm d}{\mathrm dx}$ twice to $\frac1{1-x}$ and get $\sum n^2x^n=\frac{x(1+x)}{(1-x)^3}$, which is $6$ at $x=\frac12$.)
\end{solution}

\begin{solution}{ex:bertrand}
\idea These terms decrease too slowly for the ratio test (the ratio tends to $1$) and are not comparable to a Riemann series. But they have explicit antiderivatives ($\ln\ln t$ and $-\frac1{\ln t}$), so compare with integrals.
\writeup Both $t\mapsto\frac1{t\ln t}$ and $t\mapsto\frac1{t\ln^2t}$ are positive and decreasing on $[2,\infty[$, with antiderivatives $\ln\ln t\to\infty$ and
$-\frac1{\ln t}\to0$. By Proposition~\ref{prop:posseries}(c), the first series diverges and the second converges.
\end{solution}

\begin{solution}{ex:ln2}
\idea We know the asymptotics of $H_n$ precisely (Proposition~\ref{prop:gamma}), so express the alternating partial sum through harmonic numbers: all terms minus twice the even terms. Then handle odd-index partial sums separately.
\writeup Separating even and odd terms,
\[
\begin{aligned}S_{2n}&=\sum_{k=1}^{2n}\frac1k-2\sum_{k=1}^n\frac1{2k}=H_{2n}-H_n\eqbr=\ln(2n)+\gamma-\ln n-\gamma+o(1)\to\ln2.\end{aligned}
\]
Since $S_{2n+1}-S_{2n}=\frac1{2n+1}\to0$, $S_n\to\ln2$.
\end{solution}

\begin{solution}{ex:sin}
\idea (a) A recurrence $u_{n+1}=f(u_n)$ with $f(x)<x$: the sequence decreases, and its limit is a fixed point. (b) To find an equivalent of a sequence tending to $0$, look for a power $u_n^\beta$ whose increments converge to a non-zero constant; Ces\`aro then gives $u_n^\beta\sim cn$. Since $\sin x=x-\frac{x^3}6+\cdots$, the natural choice is $\beta=-2$: the cubic correction becomes a constant increment of $\frac1{u^2}$.
\writeup (a) Since $0<\sin x<x$ for $0<x\le\frac\pi2$ (Proposition~\ref{prop:sincmp}), $(u_n)$ is positive and decreasing, so it converges to some $\ell\in\left[0,\frac\pi2\right]$ with $\sin\ell=\ell$ (continuity of $\sin$). The same inequality forbids a positive fixed point, so $\ell=0$.

(b) Taylor's theorem (Theorem~\ref{thm:taylor}) gives $\sin x=x-x^3/6+O(x^5)$; hence $\sin^2x=x^2-\frac{x^4}3+O(x^6)$ as $x\to0$, so
\[
\frac1{\sin^2x}-\frac1{x^2}=\frac{x^2-\sin^2x}{x^2\sin^2x}=\frac{\frac{x^4}3+O(x^6)}{x^4+O(x^6)}\to\frac13.
\]
With $x=u_n\to0$ this gives the first claim. By the corollary of Ces\`aro, $\frac1{nu_n^2}\to\frac13$, that is $u_n\sim\sqrt{\frac3n}$.
\end{solution}

\begin{solution}{ex:kempner}
\idea Positive terms: it suffices to bound the partial sums. Group the integers by number of digits: there are few admissible $k$-digit numbers ($\approx9^k$ instead of $10^k$), and each term is at most $10^{-(k-1)}$. The ratio $\frac9{10}$ makes the grouped series geometric.
\writeup There are $8\cdot9^{k-1}$ integers with $k$ digits and no digit $9$ (the first digit has $8$ choices, the others $9$), and each is at least $10^{k-1}$.
The terms are positive, so the partial sums are bounded by
\[
\sum_{k=1}^\infty\frac{8\cdot9^{k-1}}{10^{k-1}}=\frac8{1-\frac9{10}}=80.
\]
The bound is strict because already in the one-digit block, $1+1/2+\dots+1/8<8$.
(The true value is about $22.92$.)
\end{solution}

\begin{solution}{ex:rearr}
\idea As in Exercise~\ref{ex:ln2}, everything reduces to harmonic numbers. After $n$ blocks we have used $pn$ odd terms and $qn$ even terms, so the partial sum is (sum of the first $pn$ odd reciprocals) $-$ (sum of the first $qn$ even reciprocals). The $\ln$ terms do not cancel completely because $pn\ne qn$: that leftover $\frac12\ln\frac pq$ is the whole phenomenon.
\writeup Write $O_m=\sum_{k=1}^m\frac1{2k-1}=H_{2m}-\frac12H_m$ and $E_m=\sum_{k=1}^m\frac1{2k}=\frac12H_m$. By Proposition~\ref{prop:gamma},
\[
\begin{gathered}O_m=\ln2+\frac12\ln m+\frac\gamma2+o(1),\qq E_m=\frac12\ln m+\frac\gamma2+o(1).\end{gathered}
\]
After $n$ complete blocks, the partial sum is
$O_{pn}-E_{qn}=\ln2+\frac12\ln\frac{pn}{qn}+o(1)\to\ln2+\frac12\ln\frac pq$.
A partial sum that stops inside a block differs from one taken at the end of the previous block by at most $p+q$ terms, each tending to $0$, so the whole sequence of
partial sums has the same limit.
\end{solution}

\begin{solution}{ex:stirling}
\idea To show a sequence $d_n$ converges, show that $\sum(d_n-d_{n+1})$ converges absolutely. The differences are computable in closed form because $\ln(n+1)!-\ln n!=\ln(n+1)$; then integrating a finite geometric identity gives an expansion of $\ln(1+\frac1n)$ with an explicit error of order $1/n^4$, showing the differences are $O(\frac1{n^2})$. The term $\frac12\ln n$ in $d_n$ was chosen precisely to kill the $\frac1n$ term.
\writeup (a)
\[\begin{aligned}d_n-d_{n+1}&=-\ln(n+1)-\left(n+\frac12\right)\ln n\eqbr\quad+\left(n+\frac32\right)\ln(n+1)-1\\&=\left(n+\frac12\right)\ln\left(1+\frac1n\right)-1.\end{aligned}\]
By Proposition~\ref{prop:early-ftc}, integrating the finite identity
$1/(1+t)=1-t+t^2-t^3/(1+t)$ yields
\[
\ln(1+x)=x-\frac{x^2}2+\frac{x^3}3+r(x),
\qquad |r(x)|\le\frac{x^4}4\quad(0\le x\le1).
\]
Use this at $x=\frac1n$:
\[
\begin{aligned}&\left(n+\frac12\right)\left(\frac1n-\frac1{2n^2}+\frac1{3n^3}+O\left(\frac1{n^4}\right)\right)-1\\
&\quad=\left(\frac13-\frac14\right)\frac1{n^2}+O\left(\frac1{n^3}\right)\eqbr\quad=\frac1{12n^2}+O\left(\frac1{n^3}\right).\end{aligned}
\]
(b) So $\sum(d_n-d_{n+1})$ converges absolutely by comparison with $\sum\frac1{n^2}$, and $d_n=d_1-\sum_{k<n}(d_k-d_{k+1})\to d\in\R$. Exponentiating,
$\frac{n!}{\sqrt n\,(n/e)^n}=e^{d_n}\to e^d=C>0$.
\end{solution}


\begin{solution}{ex:oresme}
\idea (a) The terms $\frac1k$ decrease, so the cheapest lower bound for a block of terms is (number of terms) $\times$ (last term). Oresme's blocks double in length while their last term halves: each block weighs at least $\frac12$. (b) A term $\frac1{n^2}$ is just below a telescoping difference. (c) $H_n\approx\ln n+0.577$ is an excellent approximation, and $e^{9.42}$ is a big number.
\writeup (a) For $j\ge1$ the block $2^{j-1}<k\le2^j$ has $2^{j-1}$ terms, each at least $\frac1{2^j}$, so its sum is at least $\frac12$. Hence
$H_{2^m}=1+\sum_{j=1}^m\sum_{2^{j-1}<k\le2^j}\frac1k\ge1+\frac m2$. The partial sums are unbounded, so the series diverges (its terms are positive, Proposition~\ref{prop:posseries}(a)).

(b) For $n\ge2$, $\frac1{n^2}<\frac1{n(n-1)}=\frac1{n-1}-\frac1n$, so $\sum_{n=2}^N\frac1{n^2}<1-\frac1N<1$ and $\sum_{n\le N}\frac1{n^2}<2$: bounded partial sums, convergence, sum at most $2$. (The sum is $\frac{\pi^2}6\approx1.645$, as the polynomial chapter recalls; Exercise~\ref{ex:eulerbasel} computes it without knowing that.)

(c) $H_n>10$ first happens when $\ln n+\gamma+o(1)>10$, so near $n=e^{10-\gamma}\approx e^{9.4228}\approx12367$. The exact answer is $n=12367$: $H_{12366}\approx9.99996$ and $H_{12367}\approx10.00004$. By comparison, $H_n>100$ needs about $e^{99.42}\approx1.5\times10^{43}$ terms. The harmonic series diverges, but it is in no hurry.
\end{solution}

\begin{solution}{ex:nested}
\idea A recursive sequence $u_{n+1}=f(u_n)$ with $f$ increasing is monotone, in the direction of its first step; it is bounded by the fixed point if it starts below it, because $f$ maps $[0,\ell]$ into itself. For $c=2$, the half-angle formula for cosine turns $\sqrt{2+2\cos x}$ into $2\cos\frac x2$, and the product of cosines is a telescoping product of sines in disguise.
\writeup The equation $\ell=\sqrt{c+\ell}$, with $\ell\ge0$, means $\ell^2-\ell-c=0$, so $\ell=\frac{1+\sqrt{1+4c}}2$. The map $f(x)=\sqrt{c+x}$ is increasing on $[0,\infty[$ with $f(\ell)=\ell$.
Induction: $0=u_0<u_1=\sqrt c<\ell$ (as $c<\ell^2$); and if $u_n<u_{n+1}<\ell$ then applying $f$ gives $u_{n+1}<u_{n+2}<\ell$. So $(u_n)$ is increasing and bounded above by $\ell$; it converges to some $m\le\ell$, and passing to the limit in $u_{n+1}^2=c+u_n$ gives $m^2=c+m$, so $m=\ell$.
For $c=2$: $\ell=2$, so $\sqrt{2+\sqrt{2+\cdots}}=2$.

Next, $u_0=0=2\cos\frac\pi2$, and if $u_n=2\cos\frac\pi{2^{n+1}}$ then $u_{n+1}=\sqrt{2+2\cos\frac\pi{2^{n+1}}}=\sqrt{4\cos^2\frac\pi{2^{n+2}}}=2\cos\frac\pi{2^{n+2}}$ (the cosine is positive). So $\frac{u_n}2=\cos\frac\pi{2^{n+1}}$.
Finally, from $\sin x=2\sin\frac x2\cos\frac x2$ repeated $n$ times, $\sin x=2^n\sin\frac x{2^n}\prod_{k=1}^n\cos\frac x{2^k}$. At $x=\frac\pi2$:
$1=2^n\sin\frac\pi{2^{n+1}}\prod_{k=1}^n\cos\frac\pi{2^{k+1}}$, and $2^n\sin\frac\pi{2^{n+1}}=\frac\pi2\cdot\frac{\sin t}t$ with $t=\frac\pi{2^{n+1}}\to0$. Hence
$\prod_{k=1}^n\frac{u_k}2\to\frac2\pi$, which is Vi\`ete's formula: the first infinite product in the history of mathematics.
\end{solution}

\begin{solution}{ex:newton}
\idea Newton's method replaces a curve by its tangent, so the error after a step should be governed by the curvature, that is, by the \emph{square} of the previous error: that is (a). Then (b) is just iterating $e_{n+1}\le e_n^2$, which gives $e_n\le e_1^{2^{n-1}}$, a double exponential. This is why calculators compute square roots this way.
\writeup (a) $x_{n+1}-\sqrt2=\frac{x_n^2+2-2\sqrt2x_n}{2x_n}=\frac{(x_n-\sqrt2)^2}{2x_n}$. Hence $x_{n+1}\ge\sqrt2$ for all $n\ge0$, so for $n\ge1$ the denominator $2x_n\ge2\sqrt2$ and the stated bound follows.

(b) Put $e_n=\frac{x_n-\sqrt2}{2\sqrt2}\ge0$ for $n\ge1$. Then (a) says $e_{n+1}\le e_n^2$, so $e_n\le e_1^{2^{n-1}}$ by induction.
Since $e_1<0.03034$ and $2\sqrt2<2.83$, a genuine upper bound is
$x_n-\sqrt2<2.83(0.03034)^{2^{n-1}}$.
The actual errors at $x_3\approx1.414215686$ and $x_4\approx1.41421356237469$ are about $2.12\times10^{-6}$ and $1.59\times10^{-12}$; at $x_5$ the error is below $9\times10^{-25}$.
For $100$ decimal places, $n=8$ suffices: $2.83(0.03034)^{128}<10^{-100}$.

(c) The tangent to $y=x^2-2$ at $x_n$ is $y=(x_n^2-2)+2x_n(x-x_n)$; it meets $y=0$ at $x=x_n-\frac{x_n^2-2}{2x_n}=\frac12\left(x_n+\frac2{x_n}\right)$. Newton's method is ``follow the tangent to the axis, repeat''; square-root iterations of this kind have ancient antecedents. The Old Babylonian tablet YBC~7289 (c.\,1800--1600~BC) records $\sqrt2\approx1.41421296$, but does not record the algorithm used to obtain it.
\end{solution}

\begin{solution}{ex:stolz}
\idea Ces\`aro's proof is a model: split the sum into the first $N$ terms (fixed, hence negligible after dividing by $b_n\to\infty$) and the rest (where the ratio is close to $\ell$). The theorem then computes limits of quotients ``$\frac\infty\infty$'' by looking at increments, exactly as L'H\^opital's rule (Chapter~\ref{ch:calculus}) does with derivatives.
\writeup Fix $\varepsilon>0$ and choose $N$ large enough that $b_N>0$ and $\left|\frac{a_{k+1}-a_k}{b_{k+1}-b_k}-\ell\right|\le\varepsilon$ for $k\ge N$. For $n>N$, summing the inequalities $|a_{k+1}-a_k-\ell(b_{k+1}-b_k)|\le\varepsilon(b_{k+1}-b_k)$ for $N\le k<n$ (the differences $b_{k+1}-b_k$ are positive) gives
$|a_n-a_N-\ell(b_n-b_N)|\le\varepsilon(b_n-b_N)\le\varepsilon b_n$. Dividing by $b_n>0$,
\[
\left|\frac{a_n}{b_n}-\ell\right|\le\frac{|a_N-\ell b_N|}{b_n}+\varepsilon,
\]
and the first term tends to $0$ as $n\to\infty$ because $b_n\to\infty$. So $\left|\frac{a_n}{b_n}-\ell\right|\le2\varepsilon$ for $n$ large.

For $a_n=1^k+\dots+n^k$ and $b_n=n^{k+1}$: $\frac{a_{n+1}-a_n}{b_{n+1}-b_n}=\frac{(n+1)^k}{(n+1)^{k+1}-n^{k+1}}$, and the denominator is $(k+1)n^k+\dots$ by the binomial formula, so the ratio tends to $\frac1{k+1}$.
For $a_n=\ln n!$ and $b_n=n\ln n$: $\frac{\ln(n+1)}{(n+1)\ln(n+1)-n\ln n}=\frac{\ln(n+1)}{\ln(n+1)+n\ln(1+\frac1n)}$, and $n\ln(1+\frac1n)\to1$ while $\ln(n+1)\to\infty$, so the ratio tends to $1$. (Stirling's formula says more: $\ln n!=n\ln n-n+O(\ln n)$.)
\end{solution}

\begin{solution}{ex:condensation}
\idea Oresme's grouping, made into a theorem: for a non-increasing sequence, the block of $2^k$ terms from index $2^k$ is squeezed between $2^ku_{2^{k+1}}$ and $2^ku_{2^k}$. So the series and its ``condensed'' version have bounded partial sums together. The condensed series is often geometric, which is why the test is so effective for logarithms: $\ln2^k=k\ln2$ turns a log into a linear factor.
\writeup For $k\ge0$, the $2^k$ terms with $2^k\le n<2^{k+1}$ satisfy $u_{2^{k+1}}\le u_n\le u_{2^k}$, so
\[
2^ku_{2^{k+1}}\le\sum_{n=2^k}^{2^{k+1}-1}u_n\le2^ku_{2^k}.
\]
Summing over $k<K$: $\frac12\sum_{k=1}^{K}2^ku_{2^k}\le\sum_{n<2^K}u_n\le\sum_{k<K}2^ku_{2^k}$. Both series have non-negative terms, so each has bounded partial sums if and only if the other does (Proposition~\ref{prop:posseries}(a)).

For $u_n=n^{-\alpha}$ ($\alpha>0$): $2^ku_{2^k}=2^{k(1-\alpha)}$, a geometric series, convergent iff $\alpha>1$.
For $u_n=\frac1{n\ln n(\ln\ln n)^\beta}$ ($n\ge3$), which is eventually non-increasing for every fixed real $\beta$: $2^ku_{2^k}=\frac1{k\ln2\,(\ln(k\ln2))^\beta}$. The resulting sequence is also eventually non-increasing. Condense again: with $v_k=\frac1{k\ln2(\ln k+\ln\ln2)^\beta}$, $2^jv_{2^j}=\frac1{\ln2\,(j\ln2+\ln\ln2)^\beta}\sim\frac{1}{(\ln2)^{1+\beta}j^\beta}$, a Riemann series, convergent iff $\beta>1$. (The intermediate $\ln\ln2<0$ changes nothing for $k$ large, and finitely many terms never matter.)
\end{solution}

\begin{solution}{ex:sinn}
\idea Convergence is Dirichlet's test, already done in the text. For divergence, one cannot say $|\sin n|\ge c>0$ (the integers $n$ come arbitrarily close to multiples of $\pi$), but $\sin^2n$ is on average $\frac12$, and the identity $\sin^2n=\frac{1-\cos2n}2$ makes this exact: a divergent harmonic piece plus a convergent oscillating piece.
\writeup The example of the section with $\theta=1$ gives the convergence of $\sum\frac{\sin n}n$ (the partial sums of $\sin n$ are bounded by $\frac1{\sin\frac12}\approx2.09$).
For the absolute values: $|\sin n|\ge\sin^2n=\frac12-\frac{\cos2n}2$, so $\sum_{n\le N}\frac{|\sin n|}n\ge\frac12H_N-\frac12\sum_{n\le N}\frac{\cos2n}n$. The last sum converges as $N\to\infty$, by Dirichlet's test: the partial sums of $\cos2n$ are bounded (multiply by $2\sin1$ and telescope, as for the sines), and $\frac1n$ decreases to $0$. Since $H_N\to\infty$, the partial sums of $\sum\frac{|\sin n|}n$ are unbounded.
\end{solution}

\begin{solution}{ex:riemannrearr}
\idea (a) is bookkeeping: $\sum u_n$ converges and $\sum|u_n|$ does not, and $p_n,q_n$ are half-sums of these. (b) is a greedy algorithm: overshoot $S$ with positive terms, undershoot with negative ones, repeat. It never gets stuck because each pile is inexhaustible, and the overshoots shrink because $u_n\to0$. The picture to keep is a drunk walking towards a lamppost, taking shorter and shorter steps.
\writeup (a) $p_n=\frac{|u_n|+u_n}2$ and $q_n=\frac{|u_n|-u_n}2$. If $\sum p_n$ converged, then $\sum q_n=\sum(p_n-u_n)$ would converge too (difference of two convergent series), and so would $\sum|u_n|=\sum(p_n+q_n)$, contrary to the hypothesis. Same for $\sum q_n$. Both have non-negative terms and unbounded partial sums, so they diverge to $+\infty$.

(b) Let $P_1,P_2,\dots$ be the positive terms of the series in their original order and $N_1,N_2,\dots$ the negative ones (if there were finitely many of either kind, the series would be absolutely convergent). Phase~1: add $P_1,P_2,\dots$ until the running total first exceeds $S$ (use no positive terms if the initial total $0$ already exceeds $S$); this happens after finitely many terms, since $\sum P_j=+\infty$. Phase~2: add $N_1,N_2,\dots$ until the total first drops below $S$; again finitely many terms, since $\sum N_j=-\infty$. Phase~3: continue with the next unused positive terms, and so on. Every $P_j$ and every $N_j$ is used exactly once, because, after the possible empty initial phase, each phase consumes at least one term of its kind and the phases alternate forever. Insert the next unused zero term after each phase while any remain: this uses all zero terms too, and does not change any sum.

Convergence: apart from a possible empty initial phase, a positive phase ends at $T$ with $S<T\le S+P$, where $P$ is its last term; a negative phase ends at $U$ with $S+N\le U<S$, where $N<0$ is its last term. As phases advance, these last terms are original terms $u_n$ with $n\to\infty$, hence tend to zero. Both sequences of phase endpoints therefore tend to $S$. Within a phase the partial sums move monotonically between its two endpoints, so they tend to $S$ too. Inserting zeros changes none of these estimates.

(c) Aim at a moving target: positive terms until the total exceeds $1$, then one negative term, then positive terms until the total exceeds $2$, then one negative term, and so on. Because $u_n\to0$, there is an index $J$ such that $|N_j|<1$ for every $j\ge J$. Choose a stage whose following negative term has index at least $J$. Every later negative term then lowers the sum by less than $1$, and it is inserted only after the running sum has been pushed above the next integer. From that stage on, the partial sums stay above $k-1$ once the target has passed $k$, and they tend to $+\infty$. An empty positive phase consumes nothing; the construction still uses every positive term, because a finite stock of them could not push the total past every integer against the remaining negative terms. Insert one unused zero after each stage, as in (b). Thus every term is used.
\end{solution}

\begin{solution}{ex:eulerbasel}
\idea Euler wanted the number, not the formula. He noticed that $\sum\frac{x^k}{k^2}$ converges very fast at $x=\frac12$ and that a functional equation links $x$ and $1-x$; at $x=\frac12$ the equation closes on itself and the slowly convergent $\zeta(2)$ is expressed through a fast series and $\ln2$. Everything rests on two tools: integrating a finite geometric sum with an explicit remainder, and the fundamental theorem of calculus.
\writeup (a) Fix $0\le x<1$ and $n\ge1$. For $0\le u\le t\le x$, integrate $\frac1{1-u}=\sum_{k=0}^{n-1}u^k+\frac{u^n}{1-u}$ from $0$ to $t$:
$-\ln(1-t)=\sum_{k=1}^n\frac{t^k}k+r_n(t)$ with $0\le r_n(t)=\int_0^t\frac{u^n}{1-u}\dd u\le\frac1{1-x}\cdot\frac{t^{n+1}}{n+1}$.
Divide by $t$ and integrate from $0$ to $x$ (the integrands are continuous on $]0,x]$ and bounded, so the integrals are ordinary limits as the lower bound tends to $0$):
\[
L(x)=\sum_{k=1}^n\frac{x^k}{k^2}+R_n(x),\qquad0\le R_n(x)\le\frac1{1-x}\cdot\frac{x^{n+1}}{(n+1)^2}.
\]
As $n\to\infty$, $R_n(x)\to0$, so $L(x)=\sum_{k\ge1}\frac{x^k}{k^2}$. For the limit at $1$: on one hand $L(x)\le\sum_{k\ge1}\frac1{k^2}=\zeta(2)$ for every $x<1$; on the other hand, for each fixed $K$, $L(x)\ge\sum_{k\le K}\frac{x^k}{k^2}\to\sum_{k\le K}\frac1{k^2}$ as $x\to1^-$. So every limit point of $L(x)$ as $x\to1^-$ lies between $\sum_{k\le K}\frac1{k^2}$ and $\zeta(2)$ for every $K$, hence equals $\zeta(2)$; and $L$ is increasing on $[0,1[$, so the limit exists and is $\zeta(2)$.

(b) By the fundamental theorem of calculus (Proposition~\ref{prop:early-ftc}), $L'(x)=\frac{-\ln(1-x)}x$ on $]0,1[$. Let $F(x)=L(x)+L(1-x)+\ln x\ln(1-x)$. Then, with the chain rule,
\[
F'(x)=\frac{-\ln(1-x)}x-\frac{-\ln x}{1-x}+\frac{\ln(1-x)}x-\frac{\ln x}{1-x}=0,
\]
so $F$ is constant on $]0,1[$ by Theorem~\ref{thm:mvt}(c). As $x\to0^+$: $L(x)\to0$ (from (a), $0\le L(x)\le x\,\zeta(2)$ for $x<1$), $L(1-x)\to\zeta(2)$ by (a), and $\ln x\ln(1-x)\sim-x\ln x\to0$. So $F\equiv\zeta(2)$.

(c) At $x=\frac12$: $2L\left(\frac12\right)+(\ln2)^2=\zeta(2)$, that is $\zeta(2)=(\ln2)^2+2\sum_{n\ge1}\frac1{n^22^n}$. The tail of the last series after $n=20$ is at most $\sum_{n>20}\frac1{21^2\,2^n}=\frac1{21^2\,2^{20}}\approx2.2\times10^{-9}$, so twenty terms give $\zeta(2)$ to about $4\times10^{-9}$: $(\ln2)^2\approx0.480453$, $2\sum_{n\le20}\frac1{n^22^n}\approx1.164481$, total $1.644934$, which is Euler's value of 1731. Twenty terms of $\sum\frac1{n^2}$ give $1.596$, wrong in the second decimal, since the tail is about $\frac1{20}$. Euler found the exact value $\frac{\pi^2}6$ four years later; the polynomial chapter tells that story.
\end{solution}

\subsection*{Notes on the challenges}

(1) The denominator recurrence and $a_j\ge1$ give $q_j\ge F_{j+1}$ by induction, starting with $q_{-1}=0$, $q_0=1$.
If $\gamma=p/q$ in lowest terms, its finite continued fraction must have at least the certified prefix, so $q\ge q_{1000}\ge F_{1001}$.
Binet's formula (Chapter~\ref{ch:sums}) gives $F_{1001}>10^{208}$.
This is a bound from the number of quotients alone. Stronger numerical bounds require their actual values: they do not follow from the length of a prefix.
Lagarias's survey in the sources gives the wider context of the still-open irrationality problem.

(2) Convergence: $\frac1{2^n-1}\le\frac2{2^n}$. For the identity, expand $\frac1{2^n-1}=\sum_{m\ge1}2^{-nm}$ and collect, for each $N$, the pairs $(n,m)$ with $nm=N$:
there are $d(N)$ of them. (The interchange of the two sums follows from Proposition~\ref{prop:safe-sums}, since all terms are nonnegative.) Erd\H{o}s's 1948
proof of irrationality uses nothing but this expansion and a careful count of the binary digits of the partial sums; it is in \emph{J.\ Indian Math.\ Soc.}\ 12 (1948), 63--66.

(3) The count of Exercise~\ref{ex:kempner} works for any digit $d$: there are at most $9^k$ admissible $k$-digit numbers (fewer if $d\ne0$ because the first
digit has only $8$ choices), each at least $10^{k-1}$, so the series is bounded by $\sum_k9^k10^{1-k}=90$. For a block of $b$ digits, cut each $k$-digit number
into $\lfloor k/b\rfloor$ blocks: at most $(10^b-1)^{k/b}\cdot10^{b}$ numbers avoid the block, and $\left(\frac{10^b-1}{10^b}\right)^{k/b}$ still decays
geometrically. The decay is slow: avoiding $42$ removes only about one number in a hundred at each pair of digits, so the series behaves like the
harmonic series until the numbers have hundreds of digits, which is why the sum (about $228.4463$, computed by Schmelzer and Baillie in 2008) is so large, and why
no computer could find it by adding terms. Baillie's 1979 value for the digit $9$ is $22.920676619\ldots$

\begin{solution}{pb:schwartz}
\idea A polynomial in $x_1$ with coefficients in the other variables can only vanish if either its top coefficient vanishes
(probability by induction) or $x_1$ hits one of its few roots.

\writeup (a) Induction on $n$. For $n=1$, $f$ has at most $d$ roots. For $n\ge2$, write $f=\sum_{i=0}^kx_1^ig_i(x_2,\dots,x_n)$ with
$g_k\ne0$; then $\deg g_k\le d-k$. Let $r=(r_1,r')$. By induction $P(g_k(r')=0)\le\frac{d-k}{|S|}$. When $g_k(r')\ne0$, the polynomial
$f(x_1,r')$ is non-zero of degree $k$, so for each such $r'$, at most $k$ values of $r_1$ (independent of $r'$) make it vanish.
So $P(f(r)=0)\le\frac{d-k}{|S|}+\frac k{|S|}$. For $f=\prod_{i=1}^d(x_1-s_i)$ with distinct $s_i\in S$ the probability is exactly
$\frac d{|S|}$.

(b) $\det X=\sum_\sigma\operatorname{sgn}\sigma\prod_iX_{i\sigma(i)}$. The term of $\sigma$ is non-zero exactly when every
$u_iv_{\sigma(i)}$ is an edge, that is when $\sigma$ is a perfect matching, and it is then the monomial $\pm\prod_ix_{i\sigma(i)}$.
Different $\sigma$ give different monomials, so nothing cancels: $\det X\ne0$ if and only if a perfect matching exists.

(c) Substitute independent uniform values of $\mathbb F_p$ for the $x_{ij}$ and compute the determinant modulo $p$. If there is no
perfect matching, $\det X$ is the zero polynomial and the answer is always $0$. If there is one, $\det X$ reduced mod $p$ is still
non-zero (its coefficients are $\pm1$) of degree $n$, and by (a) it vanishes with probability at most $\frac np$. So the test
says ``matching'' only when there is one: it can only miss matchings.

\tools Chapter~\ref{ch:polynomials} (roots of polynomials) and Chapter~\ref{ch:linear} (determinants), both in this volume.
\leads Lovász (1979) extended the test to general graphs with the Tutte matrix, and Mulmuley, Vazirani and Vazirani (1987) used
the same determinant, with an isolation lemma, to find matchings in parallel. Polynomial identity testing is the most famous
problem that is easy with randomness and not known to be easy without it.
\end{solution}

\begin{solution}{pb:nullstellensatz}
\idea Dividing $f$ by the polynomials $g_i=\prod_{s\in S_i}(x_i-s)$ leaves a remainder of low degree in each variable that agrees
with $f$ on the grid. If $f$ vanished on the grid, the remainder would be $0$, and every top-degree monomial of $f$ would be
divisible by some $x_i^{t_i+1}$.

\writeup (a) Induction on $n$. For $n=1$ a polynomial of degree $<|S_1|$ with $|S_1|$ roots is $0$. For $n\ge2$ write
$f=\sum_{i<|S_n|}f_i(x_1,\dots,x_{n-1})x_n^i$. For fixed $s'\in S_1\times\dots\times S_{n-1}$, the polynomial $f(s',x_n)$ has degree
$<|S_n|$ and $|S_n|$ roots, so all $f_i(s')=0$. By induction each $f_i$ is $0$.

(b) Replacing $S_i$ by a subset, we may assume $|S_i|=t_i+1$. Suppose $f$ vanishes on the grid. Let
$g_i=\prod_{s\in S_i}(x_i-s)=x_i^{t_i+1}-h_i(x_i)$ with $\deg h_i\le t_i$. As long as some monomial of $f$ contains $x_i^{t_i+1}$,
replace that factor by $h_i(x_i)$. Each replacement subtracts $g_i\cdot m$ for a monomial $m$, does not raise the total degree, and
lowers the degree in $x_i$; the process stops at $\bar f$ with $\deg_{x_i}\bar f\le t_i$ and
\[
f=\bar f+\sum_ig_iu_i,\qquad\deg(g_iu_i)\le\deg f.
\]
Since $g_i=0$ on the grid, $\bar f=f=0$ there, so $\bar f=0$ by (a). Now look at the part of degree $d=\deg f$: on the right it is
$\sum_ix_i^{t_i+1}\cdot(\text{the part of degree }d-t_i-1\text{ of }u_i)$, a sum of monomials each divisible by some
$x_i^{t_i+1}$. But $x_1^{t_1}\cdots x_n^{t_n}$ has degree $d$, is divisible by no $x_i^{t_i+1}$, and has a non-zero coefficient in $f$.
Contradiction.

(c) Let $f(x)=\prod_{i=1}^n\bigl((Ax)_i-b_i\bigr)$, of degree $n$. Its part of degree $n$ is $\prod_i\sum_jA_{ij}x_j$, in which the
coefficient of $x_1x_2\cdots x_n$ is $\sum_\sigma\prod_iA_{i\sigma(i)}=\operatorname{per}A\ne0$. Apply (b) with all $t_i=1$.

\tools Chapter~\ref{ch:polynomials} (division and roots), in this volume.
\leads Alon's paper (1999) derives a dozen classical theorems from (b), several of which follow here. The permanent lemma gives,
for instance, that for every $n\times n$ matrix with non-zero permanent over $\mathbb F_p$ and every $b$ there is a $0$--$1$ vector
$x$ with $Ax$ differing from $b$ in every coordinate (Alon and Tarsi's conjecture on nowhere-zero points for $p$ odd builds on it).
\end{solution}

\begin{solution}{pb:newton}
\idea Differentiating and reversing a real-rooted polynomial keeps it real-rooted (Rolle). Doing this cleverly reduces any three
consecutive coefficients to a quadratic, whose discriminant must be non-negative.

\writeup (a) Say a polynomial \emph{splits} if it is a product of real linear factors (constants included). If $Q$ splits, so does $Q'$:
between consecutive distinct roots there is a root of $Q'$ by Rolle, a root of multiplicity $r$ is a root of $Q'$ of multiplicity
$r-1$, and this accounts for $\deg Q-1$ roots. If $Q$ splits and has degree at most $N$, so does $x^NQ(1/x)$, whose non-zero roots are the
inverses of those of $Q$.

Let $Q=P^{(k-1)}=\sum_{j\ge k-1}a_j\frac{j!}{(j-k+1)!}x^{j-k+1}$, of degree $n-k+1$, and $R(x)=x^{n-k+1}Q(1/x)=\sum_{i=0}^{n-k+1}a_{i+k-1}\frac{(i+k-1)!}{i!}x^{n-k+1-i}$.
Differentiating $R$ $(n-k-1)$ times kills all terms with $i\ge3$ and leaves the split polynomial
\[
a_{k-1}(k-1)!\frac{(n-k+1)!}2x^2+a_kk!\,(n-k)!\,x+a_{k+1}\frac{(k+1)!}2(n-k-1)!.
\]
If its leading coefficient is $0$ the inequality reads $a_k^2\ge0$. Otherwise a split quadratic has a non-negative discriminant:
\[
a_k^2k!^2(n-k)!^2\ge a_{k-1}a_{k+1}(k-1)!(k+1)!(n-k+1)!(n-k-1)!,
\]
which is $a_k^2\ge a_{k-1}a_{k+1}\frac{k+1}k\cdot\frac{n-k+1}{n-k}$.

(b) $\sum_kc(n,k)x^k=x(x+1)\cdots(x+n-1)$ (the cycle decomposition, built by inserting $1,2,\dots,n$ one at a time: element $j$ starts a new
cycle or follows one of the $j-1$ previous elements), which splits. By (a), $c(n,k)^2\ge c(n,k-1)c(n,k+1)$. A positive log-concave
sequence is unimodal: the ratios $c(n,k+1)/c(n,k)$ are non-increasing, so they are $\ge1$ up to some $k$ and $\le1$ after.

\tools Chapter~\ref{ch:polynomials} (polynomials and Rolle's theorem) and Chapter~\ref{ch:counting} (permutations), both in this volume.
\leads Newton stated (a) in 1707 without proof; the proof above is in Hardy, Littlewood and Pólya's \emph{Inequalities}. Log-concavity of
combinatorial sequences is a large subject: Heilmann and Lieb proved that matching numbers are real-rooted, and Huh and
collaborators proved Rota's and Mason's log-concavity conjectures with Hodge theory (2012--2020).
\end{solution}

\section*{Sources and further reading}

\begin{itemize}\raggedright
\item K.~Knopp, \emph{Theory and Application of Infinite Series}, Dover, 1990 (translation of the 1947 German edition). The encyclopaedia of series: convergence tests, rearrangements, Cauchy products, Abel summation and much more.
\item E.~Hairer and G.~Wanner, \emph{Analysis by Its History}, Springer, 1996. Zeno, Oresme, Grandi, Bolzano, Cauchy, Weierstrass, Riemann: the whole story of this chapter, with the original texts quoted and the original pictures.
\item Aristotle, \emph{Physics}, Book~VI, Chapter~9 (239b), and Book~VIII, Chapter~8; Zeno's dichotomy and Aristotle's answer.
\item N.~Oresme, \emph{Quaestiones super geometriam Euclidis} (c.~1350), ed.\ H.~L.~L.~Busard, Brill, 1961; Question~2 contains the grouping proof of the divergence of the harmonic series.
\item B.~Bolzano, \emph{Rein analytischer Beweis des Lehrsatzes\dots}, Prague, 1817; English translation in S.~Russ, \emph{The Mathematical Works of Bernard Bolzano}, Oxford University Press, 2004. The intermediate value theorem proved from the least upper bound property.
\item A.-L.~Cauchy, \emph{Cours d'analyse de l'\'Ecole royale polytechnique}, Paris, 1821; English translation by R.~Bradley and C.~E.~Sandifer, Springer, 2009. Chapter~VI: the definition of convergence, the criterion, the alternating series test.
\item B.~Riemann, \emph{\"Uber die Darstellbarkeit einer Function durch eine trigonometrische Reihe}, Habilitationsschrift, 1854, published in Abh.\ K\"onigl.\ Ges.\ Wiss.\ G\"ottingen 13 (1867). Section~3 contains the rearrangement theorem.
\item N.~H.~Abel, \emph{Untersuchungen \"uber die Reihe $1+\frac m1x+\dots$}, J.\ reine angew.\ Math.\ 1 (1826), 311--339. Summation by parts (Lehrsatz~III) and the first rigorous treatment of a power series on its circle of convergence.
\item J.~Havil, \emph{Gamma: Exploring Euler's Constant}, Princeton University Press, 2003. Everything about $\gamma$ and the harmonic numbers, for pleasure, including Euler's 1734 paper \emph{De progressionibus harmonicis observationes} (Eneström index E43).
\item J.~C.~Lagarias, \emph{Euler's constant: Euler's work and modern developments}, Bull.\ Amer.\ Math.\ Soc.\ 50 (2013), 527--628. A long survey, including the open question of the irrationality of $\gamma$ and the continued fraction computations of Challenge~(1).
\item C.~E.~Sandifer, \emph{How Euler Did It: Estimating the Basel problem}, MAA Online, December 2003, reprinted in \emph{How Euler Did It}, MAA, 2007. Euler's 1731 paper \emph{De summatione innumerabilium progressionum} (E20) and the formula of Exercise~\ref{ex:eulerbasel}.
\item A.~J.~Kempner, \emph{A curious convergent series}, Amer.\ Math.\ Monthly 21 (1914), 48--50; R.~Baillie, \emph{Sums of reciprocals of integers missing a given digit}, Amer.\ Math.\ Monthly 86 (1979), 372--374; T.~Schmelzer and R.~Baillie, \emph{Summing a curious, slowly convergent series}, Amer.\ Math.\ Monthly 115 (2008), 525--540. Exercise~\ref{ex:kempner} and Challenge~(3).
\item P.~Erd\H{o}s, \emph{On arithmetical properties of Lambert series}, J.\ Indian Math.\ Soc.\ 12 (1948), 63--66. The irrationality of the constant of Challenge~(2).
\item J.~Stirling, \emph{Methodus Differentialis}, London, 1730 (English translation by I.~Tweddle, Springer, 2003); A.~de Moivre, \emph{Miscellanea Analytica}, London, 1730, with its Supplement. The two halves of Stirling's formula; Tweddle's introduction sorts out who found what.
\item C.~Tweedie, \emph{James Stirling: A Sketch of His Life and Works}, Clarendon Press, 1922. His Oxford scholarships and the history of the factorial approximation.
\item D.~H.~Fowler and E.~Robson, \emph{Square root approximations in Old Babylonian mathematics: YBC~7289 in context}, Historia Mathematica 25 (1998), 366--378.
\item W.~Rudin, \emph{Principles of Mathematical Analysis}, 3rd ed., McGraw-Hill, 1976, Chapters~3 and~4. Sequences, series, continuity and Riemann's rearrangement theorem, proved in full.
\item M.~Spivak, \emph{Calculus}, 4th ed., Publish or Perish, 2008, Chapters~7, 8 and~11. The intermediate and extreme value theorems from the least upper bound axiom, Rolle and the mean value theorem, with the counterexamples of this chapter.
\end{itemize}
